\documentclass[12pt]{amsart}
%\usepackage[margin=1in]{geometry}
\pagestyle{plain}

\usepackage{amsmath,amssymb}
\usepackage{enumerate}

\usepackage[shortlabels]{enumitem}
% Set spacing for all enumerate environments
\setlist[enumerate]{itemsep=6pt, topsep=6pt, parsep=0pt}

\usepackage{graphicx}

\usepackage[left=0.75in,right=0.75in,bottom=0.9in]{geometry}  
% \usepackage{mathptmx}      % use Times fonts if available on your TeX system
\usepackage[T1]{fontenc}

\usepackage[parfill]{parskip}
\renewcommand{\baselinestretch}{1}

\usepackage{fancyhdr}
\pagestyle{fancy}
\fancyhf{}
\renewcommand{\headrulewidth}{0pt}

% Header with NAME and ID on first page only
\fancypagestyle{firstpage}{
    \fancyhf{}
    \fancyhead[L]{\small NAME: \underline{\hspace{3in}}}
    \fancyhead[R]{\small ID: \underline{\hspace{2in}}}
    \renewcommand{\headrulewidth}{0pt}
}

%%%%%%%%%%%%%%%%%%%%%%%%%%%%%%%%%
% SECTION DIVIDERS
%%%%%%%%%%%%%%%%%%%%%%%%%%%%%%%%%
\newcommand{\divider}{
  \vspace{-1em} % Add space above
  \noindent
  \raisebox{-0.15ex}{\textbullet}%
  \hspace{0.5em}%
  \rule[0.5ex]{0.95\textwidth}{1.0pt}%
  \hspace{0.5em}%
  \raisebox{-0.15ex}{\textbullet}%
  \vspace{0em} % Adjust space below
}

\newcommand{\norm}[1]{\|#1\|}

\title{ESE 2030 FINAL EXAM : Spring 2026}
\author{prof-g}


\begin{document}
\thispagestyle{firstpage}

\maketitle

\begin{center}
    {\bf INSTRUCTIONS}
\end{center}


\begin{itemize}
\item 
No books, notes, or calculators. Use a writing utensil and logic. 

\item 
Please follow question instructions and fill in the bubble-sheet. 

\item 
Select one answer per problem.

\item 
Each problem has one best answer. 

\item 
In some problems, a wrong choice may be worth partial credit. 

\item 
Be careful to fill in the correct problem number. 

\item 
These questions are written carefully: if you detect an error, read closely and do your best.

\item 
If anything is unclear, use your best judgment. 

\item 
Cheating in any form will be dealt with severely. 
\end{itemize}

\newpage


{\bf PROBLEM 1:}
% WEEK = 10
% TOPICS = Operator 2-norm, geometric meaning of largest singular value, extremal characterization
% DIFFICULTY = Core
A matrix $A \in \mathbb{R}^{4 \times 3}$ has singular values $\sigma_1 = 6$, $\sigma_2 = 3$, $\sigma_3 = 2$.
%
Over all unit vectors $\mathbf{x} \in \mathbb{R}^3$, what is the largest possible value of $\|A\mathbf{x}\|$?

\begin{enumerate}[(A)]
\item $\sqrt{6^2 + 3^2 + 2^2} = 7$
\item $36$
\item $11$
\item $6$
\item Cannot be determined from the information given.
\end{enumerate}

% \textbf{Correct Answer:} D
%
% \textbf{Explanation:} The largest singular value $\sigma_1$ is
% \emph{defined} by the extremal property
% $$
% \sigma_1 \;=\; \max_{\|\mathbf{x}\| = 1} \|A\mathbf{x}\|,
% $$
% with the maximum attained at $\mathbf{x} = \mathbf{v}_1$, the first
% right singular vector. Geometrically: $A$ maps the unit sphere in
% $\mathbb{R}^3$ to an ellipsoid in $\mathbb{R}^4$ whose longest
% semi-axis has length $\sigma_1 = 6$. So the answer is $6$.
%
% This quantity is the \textbf{operator 2-norm} $\|A\|_2 = \sigma_1$.
% It is fundamentally different from the Frobenius norm
% $\|A\|_F = \sqrt{\sum \sigma_i^2}$: the operator norm is the
% \emph{worst-case stretching} of a single unit vector, while the
% Frobenius norm aggregates over all stretching directions.
%
% \textbf{Trick answer:} A (D3 correct-structure-wrong-detail: the
% student computed $\|A\|_F = \sqrt{36 + 9 + 4} = 7$ correctly, but
% Frobenius is the wrong norm for ``maximum of $\|A\mathbf{x}\|$
% over unit $\mathbf{x}$''.
%
% \textbf{Trick answer:} B (D6 shape-implies-property: confuses
% $\sigma_1$ with $\sigma_1^2$, the largest eigenvalue of $A^T A$;
% the eigenvalue of $A^T A$ tells you $\|A\mathbf{x}\|^2$, not
% $\|A\mathbf{x}\|$).
%
% \textbf{Trick answer:} C (D5 naive-prior: just adds the singular
% values, $6 + 3 + 2 = 11$; treats them as if they combine
% additively for a single input vector).
%
% \textbf{Trick answer:} E (D4 surface match: $U$ and $V$ are needed
% to find the maximizing vector $\mathbf{v}_1$, but they are not
% needed to know the maximum \emph{value} --- that depends only on
% $\sigma_1$. The singular values alone determine the operator
% norm).

\divider

{\bf PROBLEM 2:}
% WEEK = 6
% TOPICS = Pseudoinverse, minimum-norm least-squares characterization,
%          ker(A) perp, residual orthogonality
% DIFFICULTY = Deeper
%
Let $A$ be an $m \times n$ matrix, let $\mathbf{b} \in \mathbb{R}^m$, and
let $\hat{\mathbf{x}} = A^\dagger \mathbf{b}$, where $A^\dagger$ denotes
the pseudoinverse.  Which of the following best characterizes
$\hat{\mathbf{x}}$?
%
\begin{enumerate}[(A)]
\item $\hat{\mathbf{x}} = (A^T A)^{-1} A^T \mathbf{b}$.
\item $\hat{\mathbf{x}}$ is the unique minimizer of $\norm{A\mathbf{x} - \mathbf{b}}$ in $\mathrm{im}(A)^\perp$.
\item $\hat{\mathbf{x}}$ is the unique vector in $\mathbb{R}^n$ satisfying $A\hat{\mathbf{x}} = \mathbf{b}$.
\item $\hat{\mathbf{x}}$ is the unique vector in $\mathbb{R}^n$ such that $A\hat{\mathbf{x}}$ is the orthogonal projection of $\mathbf{b}$ onto $\mathrm{im}(A)$.
\item $\hat{\mathbf{x}}$ is the unique minimizer of $\norm{A\mathbf{x} - \mathbf{b}}$ in $\ker(A)^\perp$.
\end{enumerate}
%
% \textbf{Correct Answer:} E
%
% \textbf{Explanation:} The pseudoinverse $A^\dagger$ produces, for any
% $\mathbf{b}$, the unique vector $\hat{\mathbf{x}}$ that simultaneously
% (i) minimizes the residual $\norm{A\mathbf{x} - \mathbf{b}}$ and
% (ii) has minimum norm among all such minimizers.  When $A$ has a
% nontrivial kernel, the set of residual minimizers is the affine subspace
% $\hat{\mathbf{x}}_0 + \ker(A)$; the minimum-norm choice is the unique
% representative in $\ker(A)^\perp$.  Equivalently, $\hat{\mathbf{x}}$ is
% the unique minimizer of $\norm{A\mathbf{x} - \mathbf{b}}$ that lies in
% $\ker(A)^\perp$ --- and existence-and-uniqueness on this subspace is
% the cleanest characterization, valid whether or not exact solutions
% to $A\mathbf{x} = \mathbf{b}$ exist.
%
% [Geometric picture: $A$ collapses $\ker(A)$ to zero, so any component
% of $\mathbf{x}$ inside $\ker(A)$ adds norm without changing
% $A\mathbf{x}$.  The minimum-norm solution lives where there is nothing
% to throw away --- perpendicular to the kernel.  This is precisely the
% subspace $\mathrm{coim}(A) = \ker(A)^\perp$, on which $A$ acts as an
% isomorphism onto $\mathrm{im}(A)$.]
%
% \textbf{Trick answer:} C (the naive-prior answer: treats $A^\dagger$
%   as if it were $A^{-1}$.  Wrong on two counts --- there may be no
%   solution to $A\mathbf{x} = \mathbf{b}$ at all (so the residual must
%   be minimized rather than zeroed), and even when solutions exist
%   they need not be unique.)
% \textbf{Partial credit:} D (recognizes the projection geometry ---
%   $A\hat{\mathbf{x}}$ {\em is} the orthogonal projection of $\mathbf{b}$
%   onto $\mathrm{im}(A)$ --- but misses the uniqueness issue when
%   $\ker(A) \neq \{\mathbf{0}\}$.  Many $\mathbf{x}$ satisfy
%   $A\mathbf{x} = \Pi_{\mathrm{im}(A)}(\mathbf{b})$ (any two differ by
%   an element of $\ker(A)$), so this characterization is non-unique
%   in general.  The student has the right geometric intuition for what
%   $A\hat{\mathbf{x}}$ does, but missed the minimum-norm refinement
%   that picks out a single $\hat{\mathbf{x}}$ from the affine fiber.
%   Strongest partial-credit ladder rung: same correct geometry as
%   (B), missing the kernel-perp side condition.)
% \textbf{Trick answer:} B (mixes domain and codomain subspaces:
%   $\mathrm{im}(A)^\perp \subseteq \mathbb{R}^m$, the wrong space for
%   $\hat{\mathbf{x}}$, which lives in $\mathbb{R}^n$.  A vector in
%   $\mathbb{R}^m$ cannot be a minimizer over inputs $\mathbf{x}$.
%   Surface match on ``perp of im'' without tracking which space.)
% \textbf{Partial credit:} A (correct {\em formula} when $A$ has full
%   column rank, so it gives the right answer in that special case.
%   Misses two structural points: the formula is conditional on full
%   column rank (otherwise $A^T A$ is singular), and the formula
%   doesn't describe what $\hat{\mathbf{x}}$ {\em is}, only how to
%   compute it in a special regime.)
\divider

{\bf PROBLEM 3:}
% WEEK = 10
% TOPICS = SVD structure, left vs right singular vectors, domain vs codomain
% DIFFICULTY = Deeper
A matrix $A \in \mathbb{R}^{m \times n}$ has SVD $A = U\Sigma V^T$. Which statement most fundamentally distinguishes a right singular vector $\mathbf{v}_i$ from a left singular vector $\mathbf{u}_i$?

\begin{enumerate}[(A)]
\item $\mathbf{v}_i \in \mathbb{R}^n$ lives in the domain of $A$, while $\mathbf{u}_i \in \mathbb{R}^m$ lives in the codomain, and they are linked by $A\mathbf{v}_i = \sigma_i \mathbf{u}_i$.
\item $\mathbf{v}_i$ is an eigenvector of $A^TA$, while $\mathbf{u}_i$ is an eigenvector of $AA^T$.
\item $\mathbf{v}_i$ is mapped to $\sigma_i \mathbf{u}_i$ by $A$, while $\mathbf{u}_i$ is mapped to $\sigma_i \mathbf{v}_i$ by $A$.
\item $\mathbf{v}_i$ identifies a principal input direction that gets stretched by $A$, while $\mathbf{u}_i$ identifies the corresponding output semi-axis of the image ellipsoid.
\item $\mathbf{v}_i$ is the column of $V$, while $\mathbf{u}_i$ is the column of $U$.
\end{enumerate}

% \textbf{Correct Answer:} A
% \textbf{Explanation:} The most fundamental distinction between $\mathbf{v}_i$ and $\mathbf{u}_i$ is that
% they live in \emph{different vector spaces}: $\mathbf{v}_i \in \mathbb{R}^n$ is in the domain, and
% $\mathbf{u}_i \in \mathbb{R}^m$ is in the codomain. The relation $A\mathbf{v}_i = \sigma_i \mathbf{u}_i$
% is precisely the statement that $A$ sends the $i$-th right singular vector to a scaled copy of the
% $i$-th left singular vector. Everything else (eigenvector characterizations, geometric stretching
% interpretation, columns of $U$ and $V$) is a \emph{consequence} of this domain/codomain structure.
%
% Geometrically: $\{\mathbf{v}_i\}$ is the orthonormal frame of input directions on the unit sphere
% in $\mathbb{R}^n$; $A$ scales each by $\sigma_i$ and sends it to $\mathbf{u}_i$, the corresponding
% semi-axis of the output ellipsoid in $\mathbb{R}^m$. Whenever $m \neq n$, no $\mathbf{v}_i$ can
% equal any $\mathbf{u}_j$ — they live in spaces of different dimension.
%
% \textbf{Partial credit:} B (a true equivalent characterization, but it describes \emph{how to
%   compute} the singular vectors via the symmetric matrices $A^TA$ and $AA^T$ — a derivative fact,
%   not the fundamental distinction. A student who picked this has internalized the eigenvalue-
%   problem viewpoint but missed the domain/codomain picture)
% \textbf{Partial credit:} D (geometric correct-but-derivative — captures the picture of input
%   stretching to output semi-axis, but doesn't explicitly name that the two vectors live in
%   \emph{different} spaces. Stronger than (C) because it does separate input from output)
% \textbf{Trick answer:} E (D4 surface match: a true tautology that names \emph{where} the vectors
%   sit in the matrices $V$ and $U$, but says nothing about what distinguishes them mathematically.
%   It is the choice that ``looks right'' to a student pattern-matching on the formula $A = U\Sigma V^T$)
% \textbf{Trick answer:} C (D1 twin-pair to (B): swaps the asymmetric mapping. The correct relation
%   is $A\mathbf{v}_i = \sigma_i \mathbf{u}_i$ \emph{and} $A^T\mathbf{u}_i = \sigma_i \mathbf{v}_i$;
%   the operator that sends $\mathbf{u}_i$ back to $\sigma_i \mathbf{v}_i$ is $A^T$, not $A$.
%   Students who pattern-match on symmetry without checking domain/codomain fall here)

\divider

{\bf PROBLEM 4:}
% WEEK = 12
% TOPICS = Backpropagation, chain rule as composition of derivatives,
%          transpose as adjoint, derivative factorization
% DIFFICULTY = Core (conceptual)
In a feed-forward network, layer $k$ takes input $\mathbf{h}_{k-1}$ and produces $\mathbf{z}_k = W_k \mathbf{h}_{k-1} + \mathbf{b}_k$ followed by $\mathbf{h}_k = \sigma(\mathbf{z}_k)$, with $\sigma$ applied componentwise. Define $\boldsymbol{\delta}_k = \partial L / \partial \mathbf{z}_k$.
%
Backpropagation produces a recursion that expresses $\boldsymbol{\delta}_k$ in terms of $\boldsymbol{\delta}_{k+1}$. Which is the most accurate description of \emph{why} this recursion takes the form it does?

\begin{enumerate}[(A)]
\item Since the loss is computed at the output, a hidden layer's contribution decays geometrically with depth, and the recursion encodes that decay.
\item The map $\mathbf{z}_k \mapsto \mathbf{z}_{k+1}$ factors through $\mathbf{h}_k$, and its derivative is the product of the derivatives of the two pieces --- a diagonal matrix from componentwise $\sigma$ and the matrix $W_{k+1}$ from the affine step. %Pulling a gradient back through this composition multiplies by the transposes of these derivatives in reverse order.
\item The recursion is a consequence of gradient descent on the loss, in which the learning rate is replaced by $\mathrm{diag}(\sigma'(\mathbf{z}_k))$ to keep updates well-scaled across layers.
\item Backpropagation distributes the output error $\mathbf{h}_{\text{out}} - \mathbf{t}$ across layers in proportion to each layer's share of the network's total weight magnitude.
\item Each layer must be inverted to send information from the output back toward the input, so the recursion involves $W_{k+1}^{-1}$ acting on $\boldsymbol{\delta}_{k+1}$.
\end{enumerate}

% \textbf{Correct Answer:} B
%
% \textbf{Explanation:} Backpropagation is the chain rule applied to a composition. The
% map $\mathbf{z}_k \mapsto \mathbf{z}_{k+1}$ factors as
%   $\mathbf{z}_k \;\xrightarrow{\;\sigma\;}\; \mathbf{h}_k \;\xrightarrow{\;W_{k+1}\,\cdot\,+\,\mathbf{b}_{k+1}\;}\; \mathbf{z}_{k+1}.$
% Neither piece is a linear map: the first is a nonlinear pointwise map, the second is affine.
% But the chain rule operates on \emph{derivatives}, and at a given $\mathbf{z}_k$ each piece
% has a well-defined derivative (a linear map):
%   - The derivative of $\sigma$ at $\mathbf{z}_k$ is the diagonal matrix
%     $D_k = \mathrm{diag}(\sigma'(\mathbf{z}_k))$ (componentwise application has a diagonal
%     derivative).
%   - The derivative of the affine map $\mathbf{h} \mapsto W_{k+1}\mathbf{h} + \mathbf{b}_{k+1}$
%     is $W_{k+1}$ at every point (the bias contributes a constant and drops out).
% The chain rule multiplies these derivatives in forward order, giving $W_{k+1} D_k$, and
% pulling a gradient back through a linear map multiplies by its \emph{transpose}: that is
% the adjoint, not the inverse. Composing transposes reverses order, so
%   $\boldsymbol{\delta}_k \;=\; (W_{k+1} D_k)^T \, \boldsymbol{\delta}_{k+1} \;=\; D_k^T W_{k+1}^T \, \boldsymbol{\delta}_{k+1} \;=\; D_k W_{k+1}^T \, \boldsymbol{\delta}_{k+1}$
% (using $D_k^T = D_k$ since $D_k$ is diagonal). Every structural feature of the recursion ---
% the transpose, the reversed order, the diagonal $D_k$, the absence of the bias --- comes
% directly from this derivative-factorization picture.
%
% This is the linear-algebra heart of the entire course showing up in the deep-learning
% setting: forward through a linear map uses $W$, backward through it uses $W^T$ --- the same
% transpose that drives least squares ($A^T A \mathbf{x} = A^T \mathbf{b}$),
% the FTLA ($\ker(A)^\perp = \mathrm{im}(A^T)$), and the pseudoinverse.
%
% \textbf{Trick answer:} E (the naive ``to go backward, invert'' instinct. Wrong on the linear-
%   algebra: weight matrices are generally not square --- $W_{k+1} \in \mathbb{R}^{n_{k+1} \times n_k}$ ---
%   and even when they are, the chain rule produces the adjoint, not the inverse. Adjoint and
%   inverse coincide only for orthogonal matrices.)
% \textbf{Trick answer:} A (geometric decay with depth is a possible \emph{training-dynamics}
%   phenomenon (the vanishing-gradient effect), not a property of the recursion itself. The
%   recursion can shrink, preserve, or amplify $\boldsymbol{\delta}$ depending on $W_{k+1}$
%   and $\sigma'(\mathbf{z}_k)$.)
% \textbf{Trick answer:} D (a folk picture of error ``flowing back'' in proportion to weights;
%   no chain rule, and the proportional-to-weight-magnitude story is not what the actual
%   recursion does. This is the kind of plausible-sounding pseudo-explanation the chain rule
%   replaces.)
% \textbf{Trick answer:} C (confuses the recursion with the gradient-descent update step.
%   $\boldsymbol{\delta}_k$ is computed before any learning rate enters; the learning rate
%   appears later, when $\boldsymbol{\delta}_k$ is used to form $\partial L / \partial W_k$
%   and update weights.)
\divider


{\bf PROBLEM 5:}
% WEEK = 4
% TOPICS = Change of basis, matrix representation of a linear transformation,
%          same linear map in different coordinates, similarity
% DIFFICULTY = Core
%
Let $T:V\to V$ be a linear transformation on a finite-dimensional vector space,
and let $\mathcal{B}$ and $\mathcal{C}$ be two ordered bases of $V$. Let
$[T]_{\mathcal{B}}$ and $[T]_{\mathcal{C}}$ denote the matrices of the same
linear transformation $T$ in these two bases.
%
Which of the following best describes the relationship between
$[T]_{\mathcal{B}}$ and $[T]_{\mathcal{C}}$?

\begin{enumerate}[(A)]
\item They are equal, because they represent the same linear transformation.
\item They are inverses of each other, because changing basis reverses coordinates.
\item They have the same row-reduced echelon form, because row reduction does not depend on basis.
\item They are similar matrices.
\item They have the same columns, but written in different coordinate systems.
\end{enumerate}

% \textbf{Correct Answer:} D
%
% \textbf{Explanation:} The linear transformation $T$ is the same
% coordinate-free object in both bases, but its matrix representation
% depends on the choice of basis.  The two matrices are related by a
% change of basis: if $S$ is the change-of-coordinates matrix from
% $\mathcal{C}$-coordinates to $\mathcal{B}$-coordinates, then
% \[
%     [T]_{\mathcal{C}} = S^{-1}[T]_{\mathcal{B}}S.
% \]
% Thus $[T]_{\mathcal{B}}$ and $[T]_{\mathcal{C}}$ are similar matrices.
%
% \textbf{Trick answer:} A (confuses the linear transformation with its
%   coordinate representation.  The map $T$ is the same, but its matrix
%   depends on the basis.)
%
% \textbf{Trick answer:} B (confuses the coordinate-change matrices with
%   the matrices representing $T$.  The change-of-basis maps in opposite
%   directions are inverses of each other; the two matrices representing
%   $T$ are similar, not inverse.)
%
% \textbf{Trick answer:} E (the columns of a matrix representation are the
%   coordinate vectors of the images of basis vectors.  Changing basis
%   changes both the input basis vectors and the coordinate system used to
%   express their images, so the columns are not generally the same.)
%
% \textbf{Trick answer:} C (row-reduced echelon form is not the correct
%   relationship between two coordinate representations of the same
%   operator.  Row reduction changes a matrix by left multiplication with
%   elementary matrices; change of basis for an operator uses conjugation
%   by an invertible matrix.)
%
\divider





{\bf PROBLEM 6:}
% WEEK = 5 (TIES BACK TO WEEK 2)
% TOPICS = Inner products on function spaces, parity on symmetric
%          intervals, weighted inner products, orthogonality is
%          inner-product-dependent
% DIFFICULTY = Core (synthesis)
%
Consider a vector space $C^\infty([a,b])$ of smooth functions on an
interval $[a,b]$ with inner product $\langle \cdot , \cdot \rangle$. 
Under which of the following inner products (and intervals) is the vector $p=\cosh x$ orthogonal to the vector $q=\sinh x$?
%
\begin{enumerate}[(A)]
\item $\displaystyle \langle p, q \rangle = \int_{-1}^{1} p(x)\, q(x)\, dx$, on $C^\infty([-1,1])$
\item $\displaystyle \langle p, q \rangle = \int_{0}^{1} x\, p(x)\, q(x)\, dx$, on $C^\infty([0,1])$
\item $\displaystyle \langle p, q \rangle = \int_{0}^{1} p(x)\, q(x)\, dx$, on $C^\infty([0,1])$
\item $\displaystyle \langle p, q \rangle = \int_{-1}^{1} e^{x}\, p(x)\, q(x)\, dx$, on $C^\infty([-1,1])$
\item None of the above
\end{enumerate}
%
% \textbf{Correct Answer:} A
%
% \textbf{Explanation:} Each option is a genuine inner product on the
% stated function space, and orthogonality of the same pair $(\cosh, \sinh)$
% is a relation that depends on which inner product is in use.
% Computing each:
%
% (A) $\int_{-1}^{1} \cosh(x)\sinh(x)\, dx$. The integrand is odd
%     ($\cosh$ is even, $\sinh$ is odd, product is odd), and the
%     interval is symmetric about $0$, so the integral is $0$.
%     Orthogonal.
%
% (B) $\int_{0}^{1} \cosh(x)\sinh(x)\, dx = \tfrac{1}{2}\sinh^2(1)
%     \neq 0$. Not orthogonal --- the interval $[0,1]$ breaks the
%     symmetry that made (A) vanish.
%
% (C) $\int_{0}^{1} x \cosh(x)\sinh(x)\, dx > 0$, since $x \cosh(x)
%     \sinh(x) > 0$ on $(0,1]$. Not orthogonal --- the weight
%     biases toward positive contributions, and the interval is not
%     symmetric anyway.
%
% (D) $\int_{-1}^{1} e^{x} \cosh(x) \sinh(x)\, dx$. The interval is
%     symmetric, but the weight $e^{x}$ is neither even nor odd, so
%     the combined integrand $e^{x} \cosh(x) \sinh(x)$ is not odd.
%     A direct computation gives $e^{x}\cosh(x)\sinh(x) = \tfrac{1}{2}
%     e^{x}\sinh(2x)$, which evaluates to a nonzero quantity. Not
%     orthogonal --- symmetry of the interval alone does not force
%     the integral to vanish; the \emph{combined} integrand has to
%     be odd.
%
% [Pedagogical remark: students often hold the unstated mental model
% that ``$\cosh$ and $\sinh$ are orthogonal'' as if it were an
% intrinsic property of the pair. The subtlety this problem
% highlights is that orthogonality is a relation between three
% objects --- two vectors and an inner product --- and the
% function-space lesson is the same as the one Wk 5 made for
% $\mathbb{R}^n$ with the weighted dot product: change the inner
% product, and orthogonality changes. The parity argument for (A)
% is the cleanest and most honest path to seeing why it works:
% odd integrand on a symmetric interval. Without that combination,
% the integral is generically nonzero.]
%
% \textbf{Trick answer:} C (D5 naive-prior: students who pattern-match
%   on ``smooth functions, $L^2$ integral'' without checking which
%   interval pick this. The interval $[0,1]$ is a perfectly valid
%   choice for $L^2$ on $C^\infty([0,1])$, but loses the parity
%   structure that makes the integral in (A) vanish.)
% \textbf{Trick answer:} B (D7 surface-match on ``weighted inner
%   product'' --- students who recall that weighted inner products
%   appeared in Wk 5 may grab this without checking that the
%   weight $x$ on $[0,1]$ does not create orthogonality between
%   $\cosh$ and $\sinh$. Twin pair to (B): same interval, extra
%   weight that does not help.)
% \textbf{Trick answer:} D (the deepest trap. Students who learned
%   the heuristic ``odd integrand on a symmetric interval integrates
%   to zero'' may apply it incompletely: the unweighted product
%   $\cosh(x)\sinh(x)$ is odd, and the interval $[-1,1]$ is symmetric,
%   so a student who stops checking there picks (D). The miss is
%   that the \emph{weight} $e^{x}$ is asymmetric, and the combined
%   integrand $e^{x}\cosh(x)\sinh(x)$ is no longer odd. A careful
%   student decomposes $e^{x} = \cosh(x) + \sinh(x)$ and sees that
%   the integral splits into an even-on-symmetric piece (nonzero)
%   plus an odd-on-symmetric piece (zero), giving a nonzero total.
%   This catches students who memorized the heuristic without
%   tracking which functions it applies to.)
% \textbf{Trick answer:} E (the over-cautious
\divider 

{\bf PROBLEM 7:}
% WEEK = 8
% TOPICS = QR algorithm convergence, what the iterates approach, eigenvalue
%          extraction, similarity-preserving iteration
% DIFFICULTY = Deeper
%
Let $A \in \mathbb{R}^{n \times n}$ be a real symmetric matrix with distinct
nonzero eigenvalues. The QR algorithm produces the iteration
$$A_0 = A, \qquad A_k = Q_k R_k, \qquad A_{k+1} = R_k Q_k.$$
Each iterate $A_{k+1}$ is similar to $A_k$ (and hence to $A$), so the eigenvalues
are preserved at every step. To what does the sequence $A_k$ converge to as $k \to \infty$?

\begin{enumerate}[(A)]
\item $A_k$ converges to an upper-triangular matrix whose diagonal entries are the eigenvalues of $A$.
\item $A_k$ converges to a diagonal matrix whose diagonal entries are the singular values of $A$.
\item $A_k$ converges to a diagonal matrix whose diagonal entries are the eigenvalues of $A$.
\item $A_k$ converges to $A^T A$, since the QR iteration is built from products of $Q$ and $R$.
\item $A_k$ converges to $A$ itself, since the iteration is a similarity transformation and $A$ is the unique fixed point.
\end{enumerate}

% \textbf{Correct Answer:} C
%
% \textbf{Explanation:} The QR algorithm produces a sequence of matrices, each
% similar to $A$, that converges to a form displaying the eigenvalues of $A$
% directly on the diagonal. For a general matrix the limit is upper-triangular
% (real Schur form, with the eigenvalues on the diagonal). For a symmetric
% matrix --- as here --- the limit is strictly stronger: $A_k$ converges to a
% {\em diagonal} matrix whose diagonal entries are the eigenvalues of $A$ in
% (typically) descending magnitude order.
%
% Why diagonal in the symmetric case: each iterate $A_{k+1} = Q_k^T A_k Q_k$ is
% itself symmetric (a similarity by an orthogonal matrix preserves symmetry),
% so the limit of a sequence of symmetric upper-triangular matrices is both
% symmetric and upper-triangular --- forcing all off-diagonal entries to
% vanish. So the symmetric case collapses the general upper-triangular limit
% to a diagonal limit.
%
% This is the conceptual purpose of the algorithm. Q3 P20 establishes that
% similarity preserves eigenvalues at each step; the present problem
% establishes that the limit, by being diagonal, {\em displays} those
% eigenvalues. Together, the two halves explain how the QR algorithm
% extracts eigenvalues without ever computing a characteristic polynomial.
%
% \textbf{Partial credit:} A (correctly identifies that the limit displays
%   eigenvalues on its diagonal --- the conceptual heart of the convergence
%   claim --- but states ``upper-triangular'' rather than ``diagonal.''
%   Upper-triangular is the right answer for general $A$, but the symmetric
%   structure of this problem forces the off-diagonal entries to zero. A
%   student picking (C) has the right idea about what the QR algorithm
%   produces but did not exploit the symmetric hypothesis. One rung below
%   correct.)
% \textbf{Trick answer:} E (D5 naive-prior: imagines the iteratihttps://aistudio.google.com/app/libraryon as a
%   ``settling'' that returns the original matrix. The iteration {\em is} a
%   similarity transformation at each step, but it does not fix $A$ ---
%   instead, it transforms $A$ along a path through the orbit of similar
%   matrices toward a canonical representative. The fixed points of QR
%   iteration are the matrices that are already (block-)triangular, not
%   $A$ itself.)
% \textbf{Trick answer:} B (D7 convention collision: confuses singular values
%   with eigenvalues. The QR algorithm extracts {\em eigenvalues}, not
%   singular values. For symmetric positive-definite $A$ these happen to
%   coincide in absolute value, but the algorithm extracts signed eigenvalues
%   regardless of definiteness --- so for symmetric $A$ with negative
%   eigenvalues the diagonal limit will contain those negative values, not
%   their absolute values.)
% \textbf{Trick answer:} D (D4 surface match: $Q$ and $R$ appear in the
%   iteration, and $A^T A = R_0^T Q_0^T Q_0 R_0 = R_0^T R_0$ at the first
%   step, so the symbols look related --- but the iteration does not converge
%   to $A^T A$. For symmetric $A$, $A^T A = A^2$ has eigenvalues $\lambda_i^2$,
%   not $\lambda_i$, so this answer is doubly wrong: wrong limit, wrong
%   eigenvalue magnitudes.)
\divider 

{\bf PROBLEM 8:}
% WEEK = 10
% TOPICS = Best rank-k approximation, image of truncated SVD, geometric structure
% DIFFICULTY = Deeper
The matrix $A \in \mathbb{R}^{500 \times 200}$ has SVD $A = U\Sigma V^T$ with exactly $100$ nonzero singular values $\sigma_1 \geq \sigma_2 \geq \cdots \geq \sigma_{100} > 0$ and $\sigma_{101} = \cdots = \sigma_{200} = 0$.

Let $A_{50}$ denote the best rank-$50$ approximation to $A$. Which of the following best describes the image of the unit sphere in the domain under $A_{50}$?

\begin{enumerate}[(A)]
\item A solid $50$-dimensional ellipsoid in $\mathrm{span}(\mathbf{v}_{51}, \ldots, \mathbf{v}_{100}) \subset \mathbb{R}^{500}$, with semi-axis lengths $\sigma_{51}, \ldots, \sigma_{100}$.
\item A solid $50$-dimensional ellipsoid in $\mathrm{span}(\mathbf{v}_1, \ldots, \mathbf{v}_{50}) \subset \mathbb{R}^{200}$, with semi-axis lengths $\sigma_1, \ldots, \sigma_{50}$.
\item A solid $50$-dimensional ellipsoid in $\mathrm{span}(\mathbf{u}_1, \ldots, \mathbf{u}_{50}) \subset \mathbb{R}^{500}$, with semi-axis lengths $\sigma_1^2, \ldots, \sigma_{50}^2$.
\item A solid $50$-dimensional ellipsoid in $\mathrm{span}(\mathbf{u}_1, \ldots, \mathbf{u}_{50}) \subset \mathbb{R}^{500}$, with semi-axis lengths $\sigma_1, \ldots, \sigma_{50}$.
\item A solid $50$-dimensional ellipsoid in $\mathrm{span}(\mathbf{v}_1, \ldots, \mathbf{v}_{50}) \subset \mathbb{R}^{200}$, with semi-axis lengths $\sigma_1^2, \ldots, \sigma_{50}^2$.
\end{enumerate}

% \textbf{Correct Answer:} D
% \textbf{Explanation:} The truncated SVD is
% $A_{50} = \sum_{i=1}^{50} \sigma_i \mathbf{u}_i \mathbf{v}_i^T$,
% keeping the $50$ largest singular values. Its image is A solid $50$-dimensional subspace
% of the codomain — namely $\mathrm{span}(\mathbf{u}_1, \ldots, \mathbf{u}_{50})
% \subset \mathbb{R}^{500}$. On the unit sphere in the domain, the action of
% $A_{50}$ stretches the input direction $\mathbf{v}_i$ (for $i \leq 50$) by factor
% $\sigma_i$ and sends it to $\sigma_i \mathbf{u}_i$; all directions in
% $\ker(A_{50}) = \mathrm{span}(\mathbf{v}_{51}, \ldots, \mathbf{v}_{200})$ are
% collapsed to zero. The image is a genuine $50$-dimensional ellipsoid with
% semi-axes $\sigma_1, \ldots, \sigma_{50}$.
%
% Geometrically: the full $A$ produces a $100$-dim image ellipsoid in $\mathbb{R}^{500}$.
% Truncation to rank $50$ \emph{further collapses} that ellipsoid by zeroing the
% $50$ shortest \emph{nonzero} semi-axes, leaving a solid $50$-dim ellipsoid spanned by
% the longest $50$ semi-axes. The $100$ already-zero singular values play no role
% in the truncation — they were already part of the kernel of $A$.
%
% \textbf{Trick answer:} E (compound error twin to (D): both swaps the ambient
%   space (domain instead of codomain) AND uses $\sigma_i^2$ instead of $\sigma_i$.
%   A student making both errors at once lands here. Notably, $\sigma_i^2$ are the
%   eigenvalues of $A^TA$, whose eigenvectors \emph{do} live in the domain alongside
%   the $\mathbf{v}_i$ — so the two errors here are mutually reinforcing rather
%   than independent, which is why this is a tempting wrong answer for a student
%   who half-remembers ``$A^TA$ gives the singular values.'')
% \textbf{Trick answer:} B (D1 twin-pair to (A): swaps $\mathbf{u}$ and $\mathbf{v}$.
%   Right singular vectors $\mathbf{v}_i$ live in the domain $\mathbb{R}^{200}$, so
%   this answer locates the image in the wrong space entirely. Classic
%   domain-codomain confusion.)
% \textbf{Trick answer:} C (D7 convention collision: confuses singular values with
%   eigenvalues of $A^TA$. The semi-axis lengths are $\sigma_i$ themselves, not
%   $\sigma_i^2 = \lambda_i(A^TA)$. Right ambient space, wrong scaling.)
% \textbf{Trick answer:} A (D6 shape-implies-property: keeps the wrong half of the
%   nonzero singular values — the $50$ smallest nonzero ones rather than the $50$
%   largest. ``Best rank-$50$'' keeps the top $50$, not the next-best $50$.)

\divider

{\bf PROBLEM 9:}
% WEEK = 2
% TOPICS = Vector space axioms, derived properties, distributivity, scalar multiplication
% DIFFICULTY = Deeper
%
Which of the following properties is NOT one of the axioms of a vector space?

\begin{enumerate}[(A)]
\item $\mathbf{u} + \mathbf{v} = \mathbf{v} + \mathbf{u}$ for every $\mathbf{u}, \mathbf{v} \in V$
\item There exists a vector $\mathbf{0} \in V$ such that $\mathbf{v} + \mathbf{0} = \mathbf{v}$ for every $\mathbf{v} \in V$
\item $1 \cdot \mathbf{v} = \mathbf{v}$ for every $\mathbf{v} \in V$
\item $c \cdot (\mathbf{u} + \mathbf{v}) = c \cdot \mathbf{u} + c \cdot \mathbf{v}$ for every scalar $c$ and $\mathbf{u}, \mathbf{v} \in V$
\item $0 \cdot \mathbf{v} = \mathbf{0}$ for every $\mathbf{v} \in V$
\end{enumerate}

% \textbf{Correct Answer:} E
%
% \textbf{Explanation:} The identity $0 \cdot \mathbf{v} = \mathbf{0}$ is a 
% \emph{theorem}, not an axiom. It is provable from the axioms via 
% distributivity over scalar addition: 
% $0 \cdot \mathbf{v} = (0 + 0) \cdot \mathbf{v} = 0 \cdot \mathbf{v} + 
% 0 \cdot \mathbf{v}$, and subtracting $0 \cdot \mathbf{v}$ from both sides 
% gives $\mathbf{0} = 0 \cdot \mathbf{v}$. The other four are axioms: 
% (B) commutativity of vector addition, (C) the scalar identity axiom, 
% (D) distributivity of scalar multiplication over vector addition, and 
% (E) existence of the additive identity (zero vector).
%
% [Teaching remark: the axioms are deliberately minimal. Recognizing which 
% properties are primitive versus which are derivable is the point of an 
% axiomatic system. Note that the scalar identity $1 \cdot \mathbf{v} = 
% \mathbf{v}$ \emph{must} be assumed as an axiom — without it, one could 
% define a degenerate ``scalar multiplication'' where $c \cdot \mathbf{v} 
% = \mathbf{0}$ for all $c, \mathbf{v}$, and most other axioms would still 
% hold.]
%
% \textbf{Trick answer:} C (D4 surface-match: looks similar in form to (A) 
%   — both are scalar-times-vector identities involving a specific scalar. 
%   Students may pattern-match on shape and pick the ``$1$ version'' as 
%   derivable from the ``$0$ version'' or vice versa. Twin pair to (A); 
%   the genuine asymmetry is that (C) is irreducible while (A) follows 
%   from distributivity.)
% \textbf{Trick answer:} A (D5 naive-prior: a student who has not 
%   internalized the axiom list might guess that commutativity is a 
%   ``computational consequence'' of structure. It is not — commutativity 
%   of vector addition is an explicit axiom.)
% \textbf{Trick answer:} D (D5 naive-prior: distributivity of scalar 
%   multiplication over vector addition is an explicit axiom; this is 
%   distinct from distributivity over scalar addition, which is a 
%   \emph{separate} axiom — students sometimes conflate the two and assume 
%   one implies the other.)
% \textbf{Trick answer:} B (D5 naive-prior: students who think of the zero 
%   vector as ``just there'' rather than as something whose existence is 
%   asserted by an axiom might pick this. Existence of the additive 
%   identity is an explicit axiom in every standard formulation.)

\divider

{\bf PROBLEM 10:}
% WEEK = 7, 8 SYNTHESIS
% TOPICS = Higher-order linear ODEs, polynomial differential operator, basis of solution space, repeated roots
%
Consider the third-order linear differential equation $p(D)x = 0$, where
$p(D) = (D - 2)^2 (D + 3)$ and $D = d/dt$. Which of the following is a basis for the solution space?

\begin{enumerate}[(A)]
\item $\{e^{2t},\ e^{2t},\ e^{-3t}\}$
\item $\{e^{2t},\ e^{-3t}\}$
\item $\{e^{2t},\ t e^{2t},\ e^{-3t}\}$
\item $\{e^{2t},\ e^{2t} \ln t,\ e^{-3t}\}$
\item $\{e^{-2t},\ t e^{-2t},\ e^{3t}\}$
\end{enumerate}

% \textbf{Correct Answer:} C
%
% \textbf{Explanation:} The characteristic polynomial $p(\lambda) = (\lambda - 2)^2 (\lambda + 3)$ has roots $\lambda = 2$ (algebraic multiplicity $2$) and $\lambda = -3$ (algebraic multiplicity $1$). Each simple root contributes one basis solution $e^{\lambda t}$; a root of multiplicity $k$ contributes $k$ basis solutions $e^{\lambda t}, t e^{\lambda t}, \ldots, t^{k-1} e^{\lambda t}$. So the basis is $\{e^{2t}, t e^{2t}, e^{-3t}\}$.
%
% [Recipe: read the multiplicities off the factorization, then for each root of multiplicity $k$, generate $k$ basis solutions by multiplying $e^{\lambda t}$ by powers of $t$ from $t^0$ to $t^{k-1}$. The total count must match the degree of $p$, which equals the dimension of the solution space.]
%
% \textbf{Partial credit:} B (correctly identifies the two distinct exponentials $e^{2t}$ and $e^{-3t}$ but fails to produce a \emph{basis} --- only two functions, not enough to span a $3$-dimensional solution space. Right framework, dimension count off by one. Captures a student who knows ``each distinct root gives an exponential'' but missed that a repeated root contributes additional independent solutions.)
% \textbf{Trick answer:} A (writes $e^{2t}$ twice for the repeated root, producing a linearly \emph{dependent} set --- $\{e^{2t}, e^{2t}, e^{-3t}\}$ has rank $2$, not $3$. Diagnoses a student who knows the count must equal the multiplicity but doesn't know how independence is achieved.)
% \textbf{Trick answer:} E (sign-flips both roots: reads the factor $(D - 2)$ as eigenvalue $-2$ and $(D + 3)$ as eigenvalue $3$. The factor $(D - \lambda)$ corresponds to root $\lambda$, not $-\lambda$.)
% \textbf{Trick answer:} D (multiplies by $\ln t$ instead of $t$ for the repeated root --- the polynomial ladder $1, t, t^2, \ldots$ comes from differentiation of $e^{\lambda t}$ and the structure of Jordan blocks, not from any logarithmic source.)

\divider


{\bf PROBLEM 11:}
% WEEK = 5
% TOPICS = QR decomposition, upper-triangular R, Gram-Schmidt, partial span structure
%
Let $A$ be an $m \times n$ matrix with linearly independent columns,
and let $A = QR$ be its QR decomposition.  Why must $R$ be upper triangular
rather than lower triangular?
%
\begin{enumerate}[(A)]
\item It is a convention; the decomposition $A = QR$ also admits a lower-triangular $R$
\item Because the $k$th column of $A$ is a linear combination of the first $k$ columns of $Q$
\item Because $R = Q^T A$ and the transpose of a triangular matrix is upper triangular
\item Because the columns of $Q$ are pairwise orthogonal
\item Because the $k$th column of $A$ is orthogonal to the first $k-1$ columns of $Q$
\end{enumerate}
%
% \textbf{Correct Answer:} B
%
% \textbf{Explanation:} The QR decomposition records exactly what
% Gram-Schmidt does to the columns of $A$.  If $\mathbf{q}_1, \ldots,
% \mathbf{q}_n$ are the orthonormalized vectors and $\mathbf{a}_1,
% \ldots, \mathbf{a}_n$ are the columns of $A$, then the partial-span
% invariance of Gram-Schmidt gives
% $\mathrm{span}\{\mathbf{a}_1, \ldots, \mathbf{a}_k\} =
% \mathrm{span}\{\mathbf{q}_1, \ldots, \mathbf{q}_k\}$ for every $k$.
% Hence $\mathbf{a}_k$ is a linear combination of $\mathbf{q}_1,
% \ldots, \mathbf{q}_k$ alone, with no contribution from
% $\mathbf{q}_{k+1}, \ldots, \mathbf{q}_n$.  Reading $A = QR$
% column-by-column, the $k$th column of $R$ has zeros below row $k$ —
% precisely the upper-triangular structure.
%
% [Conceptual link: the $(k,k)$-entry of $R$ is $\norm{\mathbf{u}_k}$,
% the norm of the unnormalized Gram-Schmidt output at step $k$.  The
% above-diagonal entries $R_{j,k}$ for $j < k$ are the projection
% coefficients $\langle \mathbf{a}_k, \mathbf{q}_j \rangle$ that
% Gram-Schmidt subtracts off.  $R$ is the bookkeeping ledger of
% Gram-Schmidt; its triangularity is the recursion's visible
% fingerprint.]
%
% \textbf{Trick answer:} D (true statement, but explains the structure
%   of $Q$, not the triangularity of $R$.  Surface match for the
%   student who reaches for any QR fact without tracking which fact
%   answers the question.)
% \textbf{Trick answer:} E (twin pair to (B) — same surface
%   (``$\mathbf{a}_k$ relates to the first $k-1$ columns of $Q$'') with
%   the relationship swapped from {\em linear combination} to {\em
%   orthogonality}.  In fact the $k$th column of $A$ is generally
%   {\em not} orthogonal to the previous columns of $Q$ — it has nonzero
%   components along them, which is exactly what Gram-Schmidt subtracts
%   off.)
% \textbf{Trick answer:} C (the formula $R = Q^T A$ is correct, but the
%   rest of the reasoning is nonsense — the transpose of an upper-
%   triangular matrix is lower-triangular, not the other way around.
%   D6 shape-implies-property fallacy: invokes a transpose identity
%   that doesn't say what the student wants it to say.)
% \textbf{Trick answer:} A (the naive-prior answer for the student who
%   has not connected the triangularity of $R$ to anything structural
%   and assumes it is arbitrary.  In fact the partial-span structure
%   of Gram-Schmidt {\em forces} $R$ upper-triangular; a lower-
%   triangular $R$ would correspond to running Gram-Schmidt {\em
%   backwards} on the columns of $A$, which is a different decomposition.)
%
\divider

{\bf PROBLEM 12:}
% WEEK = 10, 11 SYNTHESIS
% TOPICS = PCA via SVD of centered data, equivalence of two computational
%          routes (eigendecomposition of covariance vs. SVD of X),
%          identification of principal components with right singular vectors
% DIFFICULTY = Deeper
%
Let $X \in \mathbb{R}^{n \times d}$ be a centered data matrix ($n$
observations, $d$ features, each column with sample mean zero). The
principal components of the data may be computed in two ways: (i) by
diagonalizing the covariance matrix $C$, or (ii) by
computing the SVD $X = U\Sigma V^T$. Suppose the SVD has singular values
$\sigma_1 \geq \sigma_2 \geq \cdots \geq \sigma_d \geq 0$. Which of the
following best describes the relationship between the two routes?

\begin{enumerate}[(A)]
\item The columns of $V$ are the principal components, and the eigenvalues of $C$ are $\sigma_i^2 / n$, because the SVD diagonalizes $X^T X$ rather than $C$ itself.
\item The two routes are unrelated computational procedures that happen to give the same numerical answer; the agreement is coincidental.
\item The columns of $V$ are the principal components, and the eigenvalues of $C$ are $\sigma_i$, since the SVD computes the same diagonalization as the covariance route.
\item The columns of $V$ are the principal components, and the eigenvalues of $C$ are $\sigma_i^2$.
\item The columns of $U$ are the principal components, and the eigenvalues of $C$ are the singular values $\sigma_i$ themselves.
\end{enumerate}

% \textbf{Correct Answer:} A
%
% \textbf{Explanation:} The SVD route and the covariance-eigendecomposition
% route are not just numerically equivalent --- they are the \emph{same
% diagonalization, viewed through two factorizations}. Computing $X^T X$
% from the SVD:
% $$X^T X = (U \Sigma V^T)^T (U \Sigma V^T) = V \Sigma^T U^T U \Sigma V^T = V (\Sigma^T \Sigma) V^T,$$
% using $U^T U = I$. The matrix $\Sigma^T \Sigma$ is the $d \times d$
% diagonal matrix with entries $\sigma_i^2$. So $X^T X = V \Lambda V^T$ is
% an orthogonal diagonalization of $X^T X$ with eigenvalues $\sigma_i^2$
% and eigenvectors $\mathbf{v}_i$ (the columns of $V$).
%
% The covariance matrix is $C = \frac{1}{n} X^T X$, so the eigenvectors
% of $C$ are the same $\mathbf{v}_i$, but the eigenvalues are scaled by
% $1/n$:
% $$C = \frac{1}{n} V (\Sigma^T \Sigma) V^T \;=\; V \Lambda_C V^T,
%   \qquad \Lambda_C = \mathrm{diag}\!\left(\tfrac{\sigma_1^2}{n}, \ldots, \tfrac{\sigma_d^2}{n}\right).$$
% Hence: principal components are the columns of $V$, the right singular
% vectors of $X$; and the eigenvalues of $C$ (the variances along principal
% directions) are $\sigma_i^2 / n$.
%
% [Two structural takeaways. (i) The principal components live in the
% \emph{feature space} $\mathbb{R}^d$ --- they are directions among
% features, not among observations --- which is exactly where the columns
% of $V$ live. The columns of $U$ live in the observation space
% $\mathbb{R}^n$ and are a different (related but distinct) object: they
% give the projections of the observations onto the principal directions,
% scaled by $\sigma_i$. (ii) The square-root relationship between
% singular values and covariance eigenvalues, $\sigma_i^2 = n \cdot
% \lambda_i(C)$, is the single most important computational fact in PCA.
% In practice one almost always computes PCA via the SVD of $X$ rather
% than forming $X^T X$ explicitly, because forming $X^T X$ squares the
% condition number and amplifies numerical error. The two routes give
% identical principal components but the SVD route is dramatically more
% stable.]
%
% \textbf{Trick answer:} B (D5 naive-prior: misses the structural
%   identity entirely. The two procedures are not coincidentally
%   compatible --- they are algebraically the same diagonalization.
%   A student who picks (A) has not seen the $X^T X = V \Sigma^T \Sigma V^T$
%   computation.)
% \textbf{Trick answer:} E (D6 shape-implies-property: swaps $U$ and
%   $V$. The columns of $U$ live in the observation space $\mathbb{R}^n$,
%   not the feature space $\mathbb{R}^d$ where principal components must
%   live. Classic feature-vs-observation confusion, parallel to the
%   $X^T X$ vs $X X^T$ trap in P32. Compounded by also dropping the
%   square: claims $\sigma_i$ are eigenvalues of $C$ rather than
%   $\sigma_i^2/n$.)
% \textbf{Partial credit:} D (gets the columns-of-$V$ identification
%   correct --- the most important structural insight --- and the
%   square correctly --- the second most important structural insight
%   --- but drops the $1/n$ normalization. A student who picked (C)
%   was thinking about $X^T X$ rather than $C = \frac{1}{n} X^T X$;
%   the answer would be exactly right if the question had asked about
%   eigenvalues of $X^T X$ instead of $C$. Two of three structural
%   facts captured; one rung below (E).)
% \textbf{Partial credit:} C (gets the columns-of-$V$ identification
%   correct, but drops both the square and the $1/n$ normalization.
%   The student recognizes that PCA and SVD are the same diagonalization
%   --- a real conceptual win --- but has not tracked how the eigenvalues
%   transform under the SVD substitution. One rung below (C); same
%   ladder.)
\divider

{\bf PROBLEM 13:}
% WEEK = 1
% TOPICS = LU decomposition, computational re-use, triangular solves
% DIFFICULTY = Core
Suppose a square matrix $A$ has been factored as $A = LU$. Which statement best describes the principal practical advantage of having this factorization in hand?

\begin{enumerate}[(A)]
\item $A = LU$ removes the need for row pivoting.
\item $A = LU$ provides $A^{-1}$ directly as $U^{-1}L^{-1}$.
\item Solving $A\mathbf{x} = \mathbf{b}$ via $A = LU$ costs fewer arithmetic operations than Gaussian elimination on the augmented matrix $[\,A \mid \mathbf{b}\,]$.
\item Once $A = LU$ is computed, additional systems $A\mathbf{x} = \mathbf{b}$ are solved by forward substitution followed by back substitution.
\item The diagonal of $U$ reveals whether $A$ is invertible.
\end{enumerate}

% \textbf{Correct Answer:} D
% \textbf{Explanation:} The signature use of an LU factorization is to amortize the
% cost of factoring across many right-hand sides. The factorization itself costs
% $O(n^3)$, the same as a single Gaussian elimination. The payoff arrives when
% one wants to solve $A\mathbf{x} = \mathbf{b}_1, A\mathbf{x} = \mathbf{b}_2, \ldots$:
% each subsequent solve is just \textbf{$L\mathbf{y} = \mathbf{b}$ followed by
% $U\mathbf{x} = \mathbf{y}$}, two triangular systems at $O(n^2)$ each. This is
% the structural insight behind LU --- triangular systems are easy, so split the
% hard work of factoring from the easy work of solving.
%
% \textbf{Trick answer:} C (captures the intuition that LU saves work, but
% misattributes the source: a single solve via LU is the \emph{same} cost as
% a single Gaussian elimination; the savings come only when the factorization
% is re-used across multiple right-hand sides.)
% \textbf{Trick answer:} E (true consequence --- since $L$ has unit diagonal,
% $\det(A) = \det(U) = \prod u_{ii}$, so $A$ is invertible iff $U$ has all nonzero
% diagonal entries --- but this is a diagnostic byproduct, not the reason one
% computes the factorization.)
% \textbf{Trick answer:} B (technically correct identity $A^{-1} = U^{-1}L^{-1}$,
% but in practice one never \emph{forms} $A^{-1}$; one solves with triangular
% substitutions. The whole point of LU is to avoid computing inverses.)
% \textbf{Trick answer:} A (the opposite is true --- pivoting is precisely why
% the more general $PA = LU$ form exists. LU does not eliminate the need for
% pivoting; it accommodates it.)
\divider

{\bf PROBLEM 14:}
% WEEK = 3
% TOPICS = Quotient spaces, equivalence classes, cosets, dimension
% DIFFICULTY = Deeper
%
Let $V$ be a finite-dimensional vector space and $W \subseteq V$ a nontrivial subspace ($W\neq V$ and $W\neq\{0\}$). 
Which of the following best describes an element of the quotient space $V/W$?

\begin{enumerate}[(A)]
\item The orthogonal projection of a vector in $V$ onto the complement of $W$
\item A set of vectors in $V$ that all differ from one another by an element of $W$
\item A vector in $V$ paired with a vector in $W$
\item A vector in $V$ that is not in $W$
\item A vector in $V$ together with a choice of representative for $W$
\end{enumerate}

% \textbf{Correct Answer:} B
%
% \textbf{Explanation:} An element of $V/W$ is a \emph{coset} (equivalence 
% class) $[\mathbf{v}] = \{\mathbf{v} + \mathbf{w} : \mathbf{w} \in W\}$. 
% Two vectors $\mathbf{u}, \mathbf{v} \in V$ represent the same element of 
% $V/W$ if and only if $\mathbf{u} - \mathbf{v} \in W$. So an element of 
% $V/W$ is genuinely a \emph{set} of vectors in $V$, all pairwise differing 
% by elements of $W$ — not a single vector, not a pair, not a projection.
%
% [Teaching remark: this is the structural shift students most often miss. 
% The quotient construction does not pick out a ``preferred representative''; 
% it identifies all vectors that differ by something in $W$. The zero element 
% of $V/W$ is the entire subspace $W$ itself, viewed as one coset. Geometrically, 
% in $V = \mathbb{R}^3$ with $W$ a line through the origin, $V/W$ is a 
% 2-dimensional space whose ``points'' are the parallel lines (cosets) of $W$.]
%
% \textbf{Trick answer:} D (captures the intuition that ``quotienting kills 
%   $W$,'' and a non-$W$ representative does specify a nonzero coset — but 
%   conflates a coset with a single representative and misses that vectors 
%   inside $W$ also represent an element of $V/W$, namely the zero coset.)
% \textbf{Trick answer:} C (D4 surface-match: pairs sound like quotient-related 
%   structure, perhaps confused with direct sum $V \oplus W$ or product 
%   construction. The quotient is a single object derived from $V$ and $W$, 
%   not a pair.)
% \textbf{Trick answer:} E (D4 surface-match: ``representative'' is the 
%   right vocabulary for the wrong question — picking a representative 
%   describes how one \emph{names} a coset, not what a coset \emph{is}. 
%   Twin pair to (A) in spirit: both confuse representative with coset.)
% \textbf{Trick answer:} A (D7 convention collision: imports an inner-product 
%   notion into the abstract quotient construction. Even when $V$ has an 
%   inner product, $V/W$ and $W^\perp$ are isomorphic but not equal — and 
%   this problem makes no inner-product assumption. Premature Week-6 import.)

\divider

{\bf PROBLEM 15:}
% WEEK = 1, 6 SYNTHESIS
% TOPICS = Solvability of Ax = b, consistency, fundamental subspaces,
%          orthogonal FTLA, ker(A^T) characterization
% DIFFICULTY = Core
%
Let $A \in \mathbb{R}^{m \times n}$ and $\mathbf{b} \in \mathbb{R}^m$. Which
of the following is equivalent to ``the system $A\mathbf{x} = \mathbf{b}$
has at least one solution''?

\begin{enumerate}[(A)]
\item All three of (A), (B), (C) are equivalent to solvability
\item The augmented matrix $[\,A \mid \mathbf{b}\,]$ has the same rank as $A$
\item $\mathbf{b} \in \mathrm{im}(A)$
\item $\mathbf{b} \perp \ker(A^T)$
\item Only (A) and (C) are equivalent; (B) is a separate, stronger condition
\end{enumerate}

% \textbf{Correct Answer:} A
%
% \textbf{Explanation:} Three faces of the same fact:
%
% (A) Definition: $A\mathbf{x} = \mathbf{b}$ has a solution iff $\mathbf{b}$
% lies in the image of $A$.
%
% (B) Orthogonal FTLA: $\mathrm{im}(A)^\perp = \ker(A^T)$, so $\mathbf{b}
% \in \mathrm{im}(A)$ iff $\mathbf{b} \perp \ker(A^T)$. This is the
% \emph{Fredholm alternative}: a linear system is solvable iff $\mathbf{b}$
% is orthogonal to every solution of the homogeneous transposed system.
%
% (C) Rank criterion: appending $\mathbf{b}$ as an extra column to $A$
% leaves the rank unchanged iff $\mathbf{b}$ is already in the span of the
% columns of $A$, i.e., iff $\mathbf{b} \in \mathrm{im}(A)$.
%
% So all three statements are equivalent to solvability, and (D) is correct.
% This is the FTLA spine appearing at the very start of the exam: the
% computational view (rank), the algebraic view (image), and the geometric
% view (orthogonality) of the same condition.
%
% [Pedagogical remark: (B) is the deepest of the three but the easiest to
% misread as a stronger condition. Students often think ``orthogonal to
% everything in $\ker(A^T)$'' is more demanding than ``in $\mathrm{im}(A)$''
% — but the orthogonal FTLA says these are exactly the same condition.]
%
% \textbf{Trick answer:} C (recognizes the definitional characterization
%   only — the most basic of the three views. Two rungs below (D).)
% \textbf{Trick answer:} B (recognizes the rank-based computational
%   characterization — the workhorse for actually checking solvability by
%   row reduction. Two rungs below (D); twin pair to (A).)
% \textbf{Trick answer:} D (recognizes the orthogonal-FTLA characterization
%   — the deepest of the three views. One rung below (D); a student who
%   picked (B) over (D) has identified the most informative single condition
%   but missed that (A) and (C) are equivalent restatements.)
% \textbf{Trick answer:} E (D5 naive-prior: treats orthogonality conditions
%   as ``stronger'' than image conditions because of the universal quantifier
%   in ``$\mathbf{b} \perp \ker(A^T)$ for every vector in $\ker(A^T)$.''
%   The orthogonal FTLA is precisely the theorem that collapses this
%   apparent gap.)

\divider

{\bf PROBLEM 16:}
% WEEK = 5
% TOPICS = Cosine similarity, scale invariance, normalization, inner product geometry
%
For nonzero vectors $\mathbf{u}, \mathbf{v} \in \mathbb{R}^n$, define
$\mathrm{sim}(\mathbf{u}, \mathbf{v}) \;=\; \displaystyle \frac{\langle \mathbf{u}, \mathbf{v} \rangle}{\|{\mathbf{u}}\|\,\|{\mathbf{v}}\|}$.

Which of the following operations is guaranteed to leave $\mathrm{sim}(\mathbf{u},\mathbf{v})$ unchanged?
%
\begin{enumerate}[(A)]
\item Replacing $\mathbf{u}$ with $5\mathbf{u}$
\item Replacing $\mathbf{u}$ with $-\mathbf{u}$
\item Replacing both $\mathbf{u}$ and $\mathbf{v}$ with $\mathbf{u} + \mathbf{c}$ and $\mathbf{v} + \mathbf{c}$ for a fixed $\mathbf{c} \in \mathbb{R}^n$
\item All of the above
\item None of the above
\end{enumerate}
%
% \textbf{Correct Answer:} A
%
% \textbf{Explanation:} Cosine similarity depends only on the {\em
% directions} of the input vectors, not their magnitudes — but it is
% sensitive to the sign of each vector and to the choice of origin.
%
% (A) For any positive scalar $c > 0$:
% $\mathrm{sim}(c\mathbf{u}, \mathbf{v}) = \frac{c \langle \mathbf{u},
% \mathbf{v}\rangle}{(c\|{\mathbf{u}}\|)\|{\mathbf{v}}\|} =
% \mathrm{sim}(\mathbf{u},\mathbf{v})$.  The factors of $c$ cancel.
%
% (B) Negation flips the sign: $\mathrm{sim}(-\mathbf{u}, \mathbf{v}) =
% -\mathrm{sim}(\mathbf{u}, \mathbf{v})$.  The numerator picks up a
% sign; the denominator does not (norm is positive).  Cosine similarity
% is invariant under {\em positive} scaling, not arbitrary scaling.
%
% (C) Translation by a common vector $\mathbf{c}$ generally changes the
% similarity.  The inner product $\langle \mathbf{u}+\mathbf{c}, \mathbf{v}+\mathbf{c}\rangle
% = \langle \mathbf{u},\mathbf{v}\rangle + \langle \mathbf{u},\mathbf{c}\rangle
% + \langle \mathbf{c},\mathbf{v}\rangle + \|{\mathbf{c}}^2$ does not
% simplify, and the norms in the denominator change too.  This is why
% PCA centers data {\em before} computing similarities — translation is
% not an invariance of cosine similarity, only of {\em centered} cosine
% similarity (a.k.a. Pearson correlation).
%
% [Geometric picture: cosine similarity = cosine of the angle the two
% vectors make at the origin.  Sliding both vectors by $\mathbf{c}$
% moves them away from the origin, which generally changes the angle
% they subtend at the origin even though their relative offset is
% unchanged.  Negation rotates a vector by $180^\circ$, flipping the
% sign of the cosine.]
%
% \textbf{Trick answer:} B (negation: the naive "scaling invariance"
%   answer that doesn't track the sign restriction.  Tempting because
%   $\|{-\mathbf{u}} = \|{\mathbf{u}}$ — students stop checking
%   after the denominator.)
% \textbf{Trick answer:} C (translation: the deepest trap.  Students who
%   think "cosine similarity is just relative geometry" import a false
%   translation invariance — a mental model that confuses cosine
%   similarity with Pearson correlation.)
% \textbf{Trick answer:} D (combines a true statement (A) with two false
%   ones; tempts students who recognize (A) and assume the "more
%   invariances the better.")
% \textbf{Trick answer:} E (the over-cautious answer; tempts students
%   who, having spotted that (B) and (C) fail, conclude that nothing is
%   safe and miss (A).)
%
\divider

{\bf PROBLEM 17:}
% WEEK = 2
% TOPICS = Direct sum decomposition, subspace sum, intersection, dimension
% DIFFICULTY = Deeper
%
Let $V$ be a finite-dimensional vector space, and let $U, W \subseteq V$ be 
subspaces satisfying $U + W = V$. Which of the following are equivalent to $V = U \oplus W$ under this hypothesis?

\begin{enumerate}[(A)]
\item $U \cap W = \{\mathbf{0}\}$
\item $\dim(U) + \dim(W) = \dim(V)$
\item Every vector in $V$ can be written as $\mathbf{u} + \mathbf{w}$ with $\mathbf{u} \in U$, $\mathbf{w} \in W$
\item $U$ and $W$ together contain a basis for $V$
\item Both (A) and (B)
\end{enumerate}

% \textbf{Correct Answer:} E
%
% \textbf{Explanation:} A direct sum $V = U \oplus W$ requires \emph{both} 
% (i) every $\mathbf{v} \in V$ decomposes as $\mathbf{u} + \mathbf{w}$, and 
% (ii) the decomposition is \emph{unique}. Given that $U + W = V$ already 
% gives existence, uniqueness is the additional requirement — and uniqueness 
% is equivalent to $U \cap W = \{\mathbf{0}\}$. In finite dimensions, this is 
% \emph{also} equivalent to the dimension count $\dim(U) + \dim(W) = \dim(V)$ 
% (since $\dim(U+W) = \dim(U) + \dim(W) - \dim(U \cap W)$). So (A) and (B) 
% are each individually equivalent to the direct-sum condition under the 
% hypothesis $U + W = V$, and both must be true.
%
% [Teaching remark: students often treat ``direct sum'' as a synonym for 
% ``sum,'' missing that the directness is precisely the uniqueness of 
% decomposition. The trivial-intersection condition and the dimension-count 
% condition are two faces of the same coin, linked by the inclusion-exclusion 
% formula for subspace dimensions. Recognizing that two conditions describe 
% the same structural fact is the FTLA-style pattern of equivalent 
% characterizations applied at the level of subspaces.]
%
% \textbf{Partial credit:} B (recognizes the dimension-count characterization 
%   — captures the right structural insight but misses that (B) is equivalent 
%   under the stated hypothesis. One ladder rung below.)
% \textbf{Partial credit:} A (recognizes the trivial-intersection 
%   characterization — captures the other equivalent insight but misses that 
%   (A) is equivalent under the stated hypothesis. Twin pair to (A); same 
%   ladder rung.)
% \textbf{Trick answer:} D (D4 surface-match: ``contains a basis'' sounds 
%   structurally relevant but is too loose — $U \cup W$ contains a basis 
%   for $V$ whenever $U + W = V$, regardless of directness. The trap is 
%   pattern-matching on basis vocabulary.)
% \textbf{Trick answer:} C (this is just a restatement of $U + W = V$, 
%   which was \emph{given} as a hypothesis — picking this means the student 
%   did not recognize that something \emph{additional} is needed for 
%   directness. D5 naive-prior: existence without uniqueness.)

\divider

{\bf PROBLEM 18:}
% WEEK = 7
% TOPICS = Spectral mapping, eigenvalues of polynomials in A, eigenvectors preserved,
%          consequences for matrix exponential and matrix powers
% DIFFICULTY = Deeper
%
Let $A \in \mathbb{R}^{n \times n}$ and let $p(t) = c_k t^k + \cdots + c_1 t
+ c_0$ be a polynomial. Suppose $A\mathbf{v} = \lambda \mathbf{v}$ for a
nonzero $\mathbf{v}$. Which of the following must be TRUE?

\begin{enumerate}[(A)]
\item $p(A)$ has the same eigenvalues as $A$, since $p$ does not change the spectrum.
\item $p(A)\mathbf{v} = p(\lambda)\,\mathbf{v}$, but the eigenvalues of $p(A)$ are unrelated to those of $A$.
\item Eigenvalues of $p(A)$ are $p(\lambda_i)$ only when $A$ is diagonalizable.
\item $p(A)\mathbf{v} = \lambda\, p(\mathbf{v})$, since polynomials commute with linear maps.
\item $p(A)\mathbf{v} = p(\lambda)\,\mathbf{v}$, and the eigenvectors of $A$ are eigenvectors of $p(A)$.
\end{enumerate}

% \textbf{Correct Answer:} E
%
% \textbf{Explanation:} The \emph{spectral mapping theorem for polynomials}:
% if $A\mathbf{v} = \lambda \mathbf{v}$, then $A^j \mathbf{v} = \lambda^j
% \mathbf{v}$ for every $j \geq 0$ (induction: $A^{j+1}\mathbf{v} = A(A^j
% \mathbf{v}) = A(\lambda^j \mathbf{v}) = \lambda^{j+1}\mathbf{v}$), so
% applying $p$ termwise gives
% $$p(A)\mathbf{v} \;=\; \sum_{j=0}^{k} c_j A^j \mathbf{v} \;=\; \sum_{j=0}^{k} c_j \lambda^j \mathbf{v} \;=\; p(\lambda)\,\mathbf{v}.$$
% So $\mathbf{v}$ is an eigenvector of $p(A)$ with eigenvalue $p(\lambda)$.
% The eigenvectors of $A$ remain eigenvectors of $p(A)$; the eigenvalues
% transform by $p$.
%
% This single structural fact underlies an enormous amount of the course:
% the matrix exponential $e^{At} = \sum t^j A^j / j!$ has the same
% eigenvectors as $A$ with eigenvalues $e^{\lambda t}$; the resolvent
% $(zI - A)^{-1}$ has eigenvalues $1/(z - \lambda)$; powers $A^k$ have
% eigenvalues $\lambda^k$. The rule ``$A$ acts on its own eigenvectors as
% multiplication by $\lambda$'' propagates through every analytic function
% of $A$ that the course encounters.
%
% [Pedagogical remark: the result holds without any diagonalizability
% hypothesis — it is a per-eigenvector statement, not a global
% diagonalization statement. (D) gets this wrong by overgeneralizing the
% diagonalizable case.]
%
% \textbf{Trick answer:} B (D3 correct-structure-wrong-detail: gets the
%   identity $p(A)\mathbf{v} = p(\lambda)\mathbf{v}$ correct but then
%   claims the spectra are unrelated. The spectra are related precisely
%   by $p$: every eigenvalue of $p(A)$ is of the form $p(\lambda)$ for
%   some eigenvalue $\lambda$ of $A$.)
% \textbf{Trick answer:} D (D7 type confusion: ``$p(\mathbf{v})$'' is
%   undefined — $p$ is a polynomial in a matrix variable, not in a
%   vector variable. The expression $\lambda \cdot p(\mathbf{v})$ has
%   no meaning. Catches students who reach for ``polynomials commute''
%   without checking what objects $p$ can be applied to.)
% \textbf{Trick answer:} C (D5 naive-prior: thinks diagonalizability is
%   needed for a per-eigenvector spectral identity. The identity holds
%   eigenvector-by-eigenvector; whether the eigenvectors span
%   $\mathbb{R}^n$ is irrelevant for this single-eigenvector claim.)
% \textbf{Trick answer:} A (D5 naive-prior at the spectrum level:
%   confuses ``preserves eigenvectors'' with ``preserves eigenvalues.''
%   $A^2$ has the same eigenvectors as $A$ but eigenvalues $\lambda^2$,
%   which generically differ from $\lambda$.)

\divider

{\bf PROBLEM 19:}
% WEEK = 6
% TOPICS = Orthogonal FTLA, fundamental subspace relations, ker/im perp characterization
%
Let $A$ be an $m \times n$ matrix.  
Suppose $\mathbf{b} \in \mathbb{R}^m$ has the property that $\mathbf{b} \perp \mathrm{im}(A)$.  
Which of the following must also be true of $\mathbf{b}$?
%
\begin{enumerate}[(A)]
\item $\mathbf{b} \in \mathrm{im}(A^T)$
\item $A^T \mathbf{b} = \mathbf{0}$
\item $A \mathbf{b} = \mathbf{0}$
\item $\mathbf{b} \in \ker(A)$
\item Both (B) and (C)
\end{enumerate}
%
% \textbf{Correct Answer:} B
%
% \textbf{Explanation:} The orthogonal version of the FTLA states
% $\mathrm{im}(A)^\perp = \ker(A^T)$. So $\mathbf{b} \perp \mathrm{im}(A)$
% means $\mathbf{b} \in \ker(A^T)$, equivalently $A^T \mathbf{b} = \mathbf{0}$.
% The condition $\mathbf{b} \perp \mathrm{im}(A)$ says $\mathbf{b}$ is
% orthogonal to every column of $A$, which is exactly the statement
% that $A^T \mathbf{b} = \mathbf{0}$.
%
% [Geometric picture: $\mathbf{b}$ lives in the codomain $\mathbb{R}^m$;
% the natural orthogonal complement of $\mathrm{im}(A) \subseteq \mathbb{R}^m$
% inside $\mathbb{R}^m$ is the kernel of the adjoint $A^T$, not anything
% involving the domain $\mathbb{R}^n$.]
%
% \textbf{Trick answer:} D ($\ker(A)$ lives in the domain $\mathbb{R}^n$,
%   not in $\mathbb{R}^m$ where $\mathbf{b}$ lives — a domain/codomain
%   confusion).
% \textbf{Trick answer:} A ($\mathrm{im}(A^T) \subseteq \mathbb{R}^n$, the
%   wrong space; this is the orthogonal complement of $\ker(A)$ in the
%   domain, not of $\mathrm{im}(A)$ in the codomain).
% \textbf{Trick answer:} C ($A\mathbf{b}$ is not even defined when
%   $\mathbf{b} \in \mathbb{R}^m$ unless $m = n$; surface match on
%   "orthogonal $\Rightarrow$ kills under multiplication" without
%   tracking which matrix to multiply by).
% \textbf{Trick answer:} E (combining a true statement with a
%   wrong-domain statement; tempts students who recognized (C) but
%   also misremembered (B) as a parallel fact).
%
\divider

{\bf PROBLEM 20:}
% WEEK = 9
% TOPICS = Spectral Theorem for symmetric matrices, real eigenvalues,
%          orthogonal diagonalization, full geometric multiplicity,
%          contrast with general real matrices
% DIFFICULTY = Deeper
Let $A \in \mathbb{R}^{n \times n}$ be a real symmetric matrix. The Spectral
Theorem asserts a package of structural facts about $A$ that can each fail
for a general (non-symmetric) real $n \times n$ matrix. Which of the
following is GUARANTEED by the Spectral Theorem for symmetric $A$?

\begin{enumerate}[(A)]
\item For every eigenvalue $\lambda$ of $A$, the geometric multiplicity of $\lambda$ equals its algebraic multiplicity.
\item All three of (A), (B), and (C) are guaranteed, and each can fail for some non-symmetric real matrix.
\item There exists an orthonormal basis of $\mathbb{R}^n$ consisting of eigenvectors of $A$.
\item Every eigenvalue of $A$ is real.
\item Only (A) and (C) are guaranteed; (B) can fail when $A$ has repeated eigenvalues.
\end{enumerate}

% \textbf{Correct Answer:} B
%
% \textbf{Explanation:} The Spectral Theorem for real symmetric matrices
% asserts that $A = Q\Lambda Q^T$ where $Q$ is orthogonal (real) and
% $\Lambda$ is diagonal (real). Reading off the three guarantees:
%
% (A) Real eigenvalues: the diagonal entries of $\Lambda$ are real.
% (B) Full geometric multiplicity: the columns of $Q$ form a basis of
%     $\mathbb{R}^n$, so the eigenspaces collectively span $\mathbb{R}^n$ ---
%     equivalently, geometric multiplicity equals algebraic multiplicity
%     for every eigenvalue (no Jordan obstruction).
% (C) Orthonormal eigenbasis: the columns of $Q$ are not just an
%     eigenbasis but an orthonormal one.
%
% Each fails for some general real matrix:
% \begin{itemize}
% \item (A) fails for the rotation $\begin{pmatrix} 0 & -1 \\ 1 & 0 \end{pmatrix}$, with eigenvalues $\pm i$.
% \item (B) fails for the Jordan block $\begin{pmatrix} 2 & 1 \\ 0 & 2 \end{pmatrix}$, with algebraic multiplicity $2$ and geometric multiplicity $1$ at $\lambda = 2$.
% \item (C) fails for $\begin{pmatrix} 1 & 1 \\ 0 & 2 \end{pmatrix}$: it has an eigenbasis $\{(1,0)^T, (1,1)^T\}$ --- so it is diagonalizable --- but no orthonormal eigenbasis.
% \end{itemize}
% All three statements are independently nontrivial guarantees, all three
% are simultaneously delivered by the Spectral Theorem, and (D) is the
% only choice that captures the full content along with the contrast.
%
% [Pedagogical note: (C) alone logically implies (A) and (B) --- if an
% orthonormal eigenbasis exists, then $A = Q\Lambda Q^T$ with real $Q$
% forces real $\Lambda$, and the existence of an eigenbasis is exactly
% full geometric multiplicity. So (C) is the strongest single statement
% on the list. The reason (D) is still the best answer is that (D) names
% the contrast with general matrices --- the structural payoff of the
% theorem --- whereas (C) alone reads like a fact about a particular $A$.
% The Spectral Theorem's content is that \emph{symmetry forces this
% package}.]
%
% \textbf{Trick answer:} C (the strongest single guarantee on the list,
%   logically implying (A) and (B); a student who picked C has identified
%   the most informative individual fact but missed the synthesis and the
%   contrast with general matrices. One rung below (D).)
% \textbf{Trick answer:} D (captures the reality of the spectrum ---
%   a genuine guarantee of the theorem, but the weakest of the three on
%   its own. Real eigenvalues alone do not yield diagonalizability, let
%   alone orthogonal diagonalizability. Two rungs below (D).)
% \textbf{Trick answer:} A (captures full geometric multiplicity ---
%   equivalent to diagonalizability, a genuine guarantee, but does not
%   address reality of the eigenvalues or orthogonality of the eigenbasis.
%   Two rungs below (D); twin pair to (A).)
% \textbf{Trick answer:} E (D5 naive-prior: imports the Jordan-block
%   intuition --- ``repeated eigenvalues can cause geometric multiplicity
%   to fall short'' --- and applies it to symmetric matrices, where the
%   Spectral Theorem precisely \emph{prevents} this. A student who picked
%   E has correctly diagnosed the general behavior tested in P19 but
%   failed to register that symmetry overrides it. The canonical Wk 8 vs.
%   Wk 9 confusion: Jordan defects happen for general matrices, never
%   for symmetric ones.)

\divider

{\bf PROBLEM 21:}
% WEEK = 7
% TOPICS = Matrix exponential, definition via power series, nilpotent matrix
% DIFFICULTY = Core
Compute $e^{-N}$ for the matrix
$N = \begin{bmatrix} 0 & 1 & 0 \\ 0 & 0 & 1 \\ 0 & 0 & 0 \end{bmatrix}.$


\bigskip

(A) $\begin{bmatrix} 1 & 1 & 1 \\ 0 & 1 & 1 \\ 0 & 0 & 1 \end{bmatrix}$
$\quad : \quad$
(B) $\begin{bmatrix} 1 & -1 & -1/2 \\ 0 & 1 & -1 \\ 0 & 0 & 1 \end{bmatrix}$
$\quad : \quad$
(C) $\begin{bmatrix} 1 & -1 & 1/2 \\ 0 & 1 & -1 \\ 0 & 0 & 1 \end{bmatrix}$


(D) $\begin{bmatrix} 1 & -1 & 1 \\ 0 & 1 & -1 \\ 0 & 0 & 1 \end{bmatrix}$
$\quad : \quad$
(E) $\begin{bmatrix} 1 & -1 & 0 \\ 0 & 1 & -1 \\ 0 & 0 & 1 \end{bmatrix}$

\bigskip

% \textbf{Correct Answer:} C
%
% \textbf{Explanation:} The matrix exponential is defined by the power series $e^A = I + A + \frac{A^2}{2!} + \frac{A^3}{3!} + \cdots$. Applied to $-N$:
% $$e^{-N} = I + (-N) + \frac{(-N)^2}{2!} + \frac{(-N)^3}{3!} + \cdots.$$
% Since $N$ is nilpotent with $N^3 = 0$, the series terminates after three terms. The key observation is that $(-N)^2 = N^2$ — squaring cancels the sign — so
% $$e^{-N} = I - N + \frac{N^2}{2}.$$
% Computing $N^2 = \begin{bmatrix} 0 & 0 & 1 \\ 0 & 0 & 0 \\ 0 & 0 & 0 \end{bmatrix}$ and assembling gives the answer in (C).
%
% \emph{Alternative path:} since $e^N$ and $e^{-N}$ commute and their product is $e^{N + (-N)} = e^0 = I$, we have $e^{-N} = (e^N)^{-1}$. A student who computes $e^N$ and inverts the resulting upper-triangular matrix arrives at the same place. This route emphasizes that $e^A$ is always invertible — a foundational fact about the matrix exponential.
%
% \textbf{Partial credit:} D (forgets the $\frac{1}{2!}$ factorial coefficient on the $N^2$ term. Knows the series terminates and the signs are right; missed the coefficient.)
% \textbf{Partial credit:} E (truncates the series at $I - N$, missing the contribution from $N^2$. Diagnoses a student who computed the first two terms and didn't check whether higher powers vanish.)
% \textbf{Trick answer:} A (twin pair to (D): computed $e^{N}$ instead of $e^{-N}$, \emph{and} forgot the factorial. The student missed the negative sign in the question entirely. No partial credit; the sign error is fundamental.)
% \textbf{Trick answer:} B (negates the sign of \emph{every} term, including the $N^2$ term, missing that $(-N)^2 = +N^2$. A subtle but fundamental error: squaring cancels signs. Catches students who reflexively flip every sign rather than thinking term-by-term.)
\divider

{\bf PROBLEM 22:}
% WEEK = 5
% TOPICS = Adjoint, definition, type signature, inner product spaces
%
Let $V$ and $W$ be inner product spaces and let $T : V \to W$ be a linear
transformation.  The adjoint $T^*$ is the linear transformation
characterized by which of the following?
%
\begin{enumerate}[(A)]
\item $T^* : V \to W$ with $T^* \circ T = I_V$
\item $T^* : W \to V$ with $T(T^*(\mathbf{w})) = \mathbf{w}$ for all $\mathbf{w} \in W$
\item $T^* : W \to V$ with $\langle T(\mathbf{v}), T^*(\mathbf{w}) \rangle = \langle \mathbf{v}, \mathbf{w} \rangle$ for all $\mathbf{v} \in V$, $\mathbf{w} \in W$
\item $T^* : W \to V$ with $\langle T(\mathbf{v}), \mathbf{w} \rangle = \langle \mathbf{v}, T^*(\mathbf{w}) \rangle$ for all $\mathbf{v} \in V$, $\mathbf{w} \in W$
\item $T^* : V \to W$ with $\langle T(\mathbf{v}), \mathbf{w} \rangle = \langle \mathbf{v}, T^*(\mathbf{w}) \rangle$ for all $\mathbf{v} \in V$, $\mathbf{w} \in W$
\end{enumerate}
%
% \textbf{Correct Answer:} D
%
% \textbf{Explanation:} The adjoint $T^*$ goes in the opposite direction
% to $T$, so $T^* : W \to V$.  Its defining property is the cross-pairing
% identity $\langle T(\mathbf{v}), \mathbf{w} \rangle_W = \langle
% \mathbf{v}, T^*(\mathbf{w}) \rangle_V$ — moving $T$ across the inner
% product produces $T^*$ on the other side.
%
% [Type-signature picture: $T(\mathbf{v}) \in W$ pairs with $\mathbf{w}
% \in W$ on the left, and $\mathbf{v} \in V$ pairs with $T^*(\mathbf{w})
% \in V$ on the right. Each inner product is taken in the space where
% both arguments live.]
%
% \textbf{Trick answer:} E (correct identity, wrong direction — twin
%   pair to (B), differing only in the type signature of $T^*$.  Catches
%   students who memorize the identity but not which space $T^*$ acts
%   between.)
% \textbf{Trick answer:} C (the {\em isometry} condition — preserving
%   inner products — not the adjoint condition.  Surface match for
%   students who confuse adjoint with the inverse of an orthogonal
%   transformation.)
% \textbf{Trick answer:} A (left-inverse property — confuses adjoint
%   with inverse, with the wrong type signature on top of it.)
% \textbf{Trick answer:} B (right-inverse property — same confusion as
%   (D) with the right type signature; tempts students who got the
%   direction right but think "adjoint = inverse on the codomain side.")
%
\divider

{\bf PROBLEM 23:}
% WEEK = 6
% TOPICS = Orthogonal complements, dimension formula, double complement, finite-dimensional inner product spaces
%
Let $V$ be a finite-dimensional inner product space and let $W \subseteq V$
be a subspace.  Which of the following must be true?
%
\begin{enumerate}[(A)]
\item $\dim(W) + \dim(W^\perp) = \dim(V)$
\item $W \cap W^\perp = \{\mathbf{0}\}$
\item All three of (A), (B), (C) must be true
\item $(W^\perp)^\perp = W$
\item Only (A) and (B) must be true; (C) requires $W$ to be closed under additional structure
\end{enumerate}
%
% \textbf{Correct Answer:} C
%
% \textbf{Explanation:} In a finite-dimensional inner product space, all
% three statements hold for every subspace $W$:
%
% (A) follows from the orthogonal decomposition $V = W \oplus W^\perp$.
%
% (B) is immediate: any $\mathbf{v} \in W \cap W^\perp$ satisfies
%     $\langle \mathbf{v}, \mathbf{v} \rangle = 0$, hence $\mathbf{v} = \mathbf{0}$
%     by positive-definiteness of the inner product.
%
% (C) follows from (A) applied twice: $\dim((W^\perp)^\perp) = \dim(V) -
%     \dim(W^\perp) = \dim(W)$, and the inclusion $W \subseteq (W^\perp)^\perp$
%     is automatic, so equality of dimensions forces equality of subspaces.
%
% [Conceptual note: $\perp$ is an involution on the lattice of subspaces
% of a finite-dimensional inner product space. The three properties
% above are not independent observations — they are the three faces of
% the orthogonal decomposition theorem.]
%
% \textbf{Partial credit:} A (correct, but the headline dimension count
%   only — misses the structural facts (B) and (C) that make the
%   complement well-behaved).
% \textbf{Partial credit:} B (correct, but the trivial-intersection fact
%   alone — recognizes the direct sum is internal but not the dimension
%   count or the involution).
% \textbf{Trick answer:} D (correct as a standalone statement, but a
%   student who picks (C) over (D) has missed (A) and (B) — the
%   double-complement identity is the deepest of the three, and the
%   pattern-matching student grabs it without checking the others).
% \textbf{Trick answer:} E (twin pair to (D) under a fabricated caveat —
%   tempts the student who half-remembers that $(W^\perp)^\perp = W$ can
%   fail in {\em infinite} dimensions and overgeneralizes the caveat to
%   the finite-dimensional setting where it does not apply).
%
\divider

{\bf PROBLEM 24:}
% WEEK = 9
% TOPICS = Perron-Frobenius theorem, dominant eigenvalue, positive eigenvector
% DIFFICULTY = Core
Suppose $A$ is a square matrix with all entries strictly positive. Which statement most completely captures what the Perron-Frobenius theorem guarantees about the eigenstructure of $A$?

\begin{enumerate}[(A)]
\item The spectral radius $\rho(A)$ is itself an eigenvalue of $A$.
\item There is a real positive eigenvalue, strictly greater in magnitude than every other eigenvalue, with a corresponding eigenvector having all strictly positive entries.
\item The largest magnitude eigenvalue is dominant.
\item There exists an eigenvector of $A$ with all strictly positive entries.
\item All eigenvalues of $A$ are real and positive.
\end{enumerate}

% \textbf{Correct Answer:} B
%
% \textbf{Explanation:} The Perron-Frobenius theorem (in its strictly-positive
% form) asserts a \textbf{bundle} of facts about the dominant eigenpair of a
% positive matrix: there is a single eigenvalue $\lambda_*$ that is
% \emph{(i)} real and positive, \emph{(ii)} equal to the spectral radius
% $\rho(A)$, \emph{(iii)} \emph{strictly} dominant (no other eigenvalue has the
% same magnitude), and \emph{(iv)} associated with an eigenvector
% $\mathbf{v}_*$ whose entries are all strictly positive. Choice (C) packages
% three of these four guarantees into a single statement; (A), (B), and (D)
% each isolate one piece and miss the others. The eigenvalue $\lambda_*$ is
% sometimes called the \emph{Perron eigenvalue} and $\mathbf{v}_*$ the
% \emph{Perron vector}.
%
% [Teaching remark: the strictness in ``strictly greater in magnitude than
% every other eigenvalue'' is what makes power iteration converge for positive
% matrices --- if the dominant eigenvalue were merely tied with another, the
% iterates would not settle into a single direction. The positivity of the
% eigenvector is what makes Perron-Frobenius the right tool for stationary
% distributions of stochastic matrices, population dynamics, and PageRank-style
% ranking applications: the entries of $\mathbf{v}_*$ can be interpreted as
% probabilities, populations, or scores.]
%
% \textbf{Trick answer:} A (captures one consequence of Perron-Frobenius
% --- that $\rho(A)$ is realized as an actual eigenvalue rather than just a
% bound --- but says nothing about strict dominance or the positive
% eigenvector. One rung of the ladder.)
% \textbf{Trick answer:} D (captures the positive-eigenvector half of the
% theorem but is silent about which eigenvalue this eigenvector corresponds
% to. A second rung of the ladder.)
% \textbf{Trick answer:} C (captures strict dominance --- the gap between
% $\rho(A)$ and the next-largest magnitude --- but says nothing about the
% sign or reality of the dominant eigenvalue, nor about the positive
% eigenvector. A third rung of the ladder.)
% \textbf{Trick answer:} E (over-generalizes ``positive matrix'' to ``all
% eigenvalues positive.'' Perron-Frobenius guarantees only that the
% \emph{dominant} eigenvalue is real and positive; the other eigenvalues can
% be complex or negative, provided their magnitudes are strictly smaller.
% For example, the $3 \times 3$ positive matrix
% $\begin{bmatrix} 1 & 1 & 1 \\ 2 & 1 & 1 \\ 1 & 2 & 1 \end{bmatrix}$ has
% one dominant real eigenvalue near $3.4$ and a complex conjugate pair with
% magnitude near $0.6$. Classic D6 shape-implies-property fallacy.)

\divider

{\bf PROBLEM 25:}
% WEEK = 11
% TOPICS = Covariance matrix, definition, statistical structure, centered data
% DIFFICULTY = Core
Let $X \in \mathbb{R}^{n \times d}$ be a centered data matrix: $n$ observations of $d$ features, with each column having mean zero. Which statement best describes the covariance matrix of $X$?

\begin{enumerate}[(A)]
\item The $d \times d$ symmetric matrix $X^T X$, whose entries are the inner products of feature columns.
\item The $d \times d$ matrix obtained by standardizing each feature to unit variance, with entry $C_{ij}$ equal to the correlation between features $i$ and $j$.
\item The $n \times n$ symmetric matrix $\frac{1}{n} X X^T$, whose entries measure pairwise relationships between observations.
\item The $d \times d$ symmetric matrix $\frac{1}{n} X^T X$, whose diagonal entries are the variances of the features and whose off-diagonal entries are the covariances between feature pairs.
\item The $n \times n$ symmetric matrix $\frac{1}{d} X X^T$, whose diagonal entries are the squared distances of each observation from the origin, normalized by the number of features.
\end{enumerate}

% \textbf{Correct Answer:} D
% \textbf{Explanation:} For a centered data matrix $X \in \mathbb{R}^{n \times d}$, the
% covariance matrix is
% $$C = \frac{1}{n} X^T X \in \mathbb{R}^{d \times d}.$$
% Its structure is most easily read off entry-by-entry: writing $X$ in column form as
% $X = [\mathbf{x}_1 \mid \mathbf{x}_2 \mid \cdots \mid \mathbf{x}_d]$ where each
% $\mathbf{x}_j \in \mathbb{R}^n$ is a centered feature column,
% $$C_{ij} = \frac{1}{n} \mathbf{x}_i^T \mathbf{x}_j = \frac{1}{n} \sum_{k=1}^{n} x_{ki} x_{kj}.$$
% The diagonal $C_{ii} = \frac{1}{n}\|\mathbf{x}_i\|^2$ is the variance of feature $i$;
% the off-diagonal $C_{ij}$ is the covariance between features $i$ and $j$.
%
% Two structural facts worth holding alongside the definition: (1) the matrix is
% $d \times d$, not $n \times n$ — it captures relationships between \emph{features},
% not between observations; (2) centering is essential — without it, $X^T X$ would
% include the means and would not give covariances at all.
%
% \textbf{Trick answer:} C (D6 shape-implies-property: swaps $X^TX$ for $XX^T$.
%   The $n \times n$ Gram-of-observations matrix is a different object — useful in
%   some kernel/dual formulations of PCA, but \emph{not} the covariance matrix.
%   A student who picks A has confused the feature-side and observation-side Gram
%   matrices, the classic dimension confusion in PCA.)
% \textbf{Trick answer:} A (D4 surface match: drops the normalization entirely.
%   $X^TX$ is proportional to the covariance matrix but is not itself the covariance
%   matrix — its entries scale with sample size. A student who picks C has lost
%   the ``per-observation'' character of variance.)
% \textbf{Trick answer:} B (D7 convention collision: describes the \emph{correlation}
%   matrix, not the covariance matrix. The two are related — correlation PCA
%   standardizes features before computing covariance — but they are distinct
%   objects with different entries and different units.)
% \textbf{Trick answer:} E (a real and nameable object — the $1/d$-normalized
%   observation Gram matrix, whose diagonal does encode squared observation
%   magnitudes — but it is the wrong-side Gram matrix \emph{and} the wrong
%   normalization. Catches a student who pattern-matches on ``$XX^T$ with some
%   normalization'' without checking which side or which scaling.)

\divider

{\bf PROBLEM 26:}
% WEEK = 10
% TOPICS = SVD block structure, four fundamental subspaces, coimage, FTLA
% DIFFICULTY = Deeper
The diagram below is an abstract representation of the SVD of a linear transformation $M$, with five regions labeled (A) through (E). Which region best represents the {\em coimage} of the transformation? (choose (A) - (E))

\begin{center}
\begin{figure}[h]
\includegraphics[width=5.5in]{SVD-ABSTRACT.png}
\end{figure}
\end{center}
\vspace{-0.1in}

% \textbf{Correct Answer:} D
% \textbf{Explanation:} Reading the diagram: $U$ is split into its first $r$ columns
% (A, shaded) and its remaining columns (B); $\Sigma$ is the rectangular middle block
% (C); and $V^T$ is split into its first $r$ rows (D, shaded) and its remaining rows
% (E), where $r = \mathrm{rank}(M)$.
%
% The coimage $\mathrm{coim}(M) = \mathbb{R}^n / \ker(M)$ is naturally represented in
% the domain by $\ker(M)^\perp$, which is spanned by the first $r$ right singular
% vectors $\mathbf{v}_1, \ldots, \mathbf{v}_r$. The rows of $V^T$ are exactly the
% $\mathbf{v}_i^T$, so the first $r$ rows of $V^T$ — region D — house the right
% singular vectors that span the coimage representative.
%
% FTLA spine: $\mathrm{coim}(M) \cong \mathrm{im}(M)$, so regions D (in the domain)
% and A (in the codomain) describe naturally isomorphic spaces. But the \emph{coimage
% itself lives in the domain}, where $V$ acts — hence D, not A.
%
% \textbf{Partial credit:} A (recognizes the rank-$r$ shaded structure and the
%   FTLA isomorphism $\mathrm{coim} \cong \mathrm{im}$, but places coimage on the
%   codomain side. A student who picked A understands the FTLA punchline but has
%   not separated the natural location of $\mathrm{coim}(M)$ from that of
%   $\mathrm{im}(M)$. This is the strongest partial-credit rung.)
% \textbf{Trick answer:} E (D1 twin-pair to (D), within the same $V^T$ block.
%   Region E is the last $n - r$ rows of $V^T$, spanning $\ker(M)$ — the orthogonal
%   complement of the coimage representative. Classic kernel-vs-coimage swap.)
% \textbf{Trick answer:} B (D6 shape-implies-property: the unshaded right portion of
%   $U$ spans $\ker(M^T) = \mathrm{im}(M)^\perp$, the natural representative of the
%   \emph{cokernel}. Confuses coimage (domain side) with cokernel (codomain side).)
% \textbf{Trick answer:} C (D4 surface match: the $\Sigma$ block holds the singular
%   values — stretching factors, not a subspace. A student who picks $C$ has read
%   ``coimage'' as ``somewhere on the diagonal'' rather than as a subspace.)

\divider

{\bf PROBLEM 27:}
% WEEK = 12
% TOPICS = Role of nonlinear activation, expressive power, why sigma is essential,
%          structural collapse without nonlinearity
% DIFFICULTY = Deeper
%
A feed-forward neural network alternates affine layers $\mathbf{z}_k = W_k \mathbf{h}_{k-1} +
\mathbf{b}_k$ with a componentwise nonlinear activation $\mathbf{h}_k = \sigma(\mathbf{z}_k)$.
Removing the activations from such a network --- so each layer is purely affine and the layers
are composed directly --- collapses the network's expressive power dramatically. Which of the
following best describes the {\em fundamental} reason a nonlinear $\sigma$ is required?

\begin{enumerate}[(A)]
\item Without $\sigma$, the Universal Approximation Theorem does not apply, so the network is not guaranteed to approximate arbitrary target functions.
\item Without $\sigma$, gradients cannot flow back through the network during backpropagation, so the weights cannot be trained.
\item Without $\sigma$, hidden activations are unbounded and the network's outputs cannot be interpreted as probabilities or classifications.
\item Without $\sigma$, the network can only fit linear targets, and most real-world relationships between inputs and outputs are nonlinear.
\item Without $\sigma$, the entire network --- regardless of depth or width --- represents only an affine function of its input, so no amount of stacking enlarges the class of functions the network can compute.
\end{enumerate}

% \textbf{Correct Answer:} E
%
% \textbf{Explanation:} The structural fact is that the composition of affine maps is
% affine. Writing two consecutive layers without activation,
% $$\mathbf{z}_2 = W_2(W_1 \mathbf{x} + \mathbf{b}_1) + \mathbf{b}_2 = (W_2 W_1) \mathbf{x}
% + (W_2 \mathbf{b}_1 + \mathbf{b}_2),$$
% the result is again of the form $W' \mathbf{x} + \mathbf{b}'$. Iterating, a network of any
% depth without $\sigma$ collapses to a single affine map $\mathbf{x} \mapsto W' \mathbf{x}
% + \mathbf{b}'$. Depth gives no new expressive power; width gives no new expressive power;
% the function class is exactly the affine maps. The nonlinear $\sigma$ is the {\em only}
% ingredient that breaks this collapse and makes depth meaningful. Every other consequence
% --- universal approximation, the ability to fit nonlinear targets, the bounded-output
% interpretation of sigmoid specifically --- is downstream of this structural fact.
%
% [Pedagogical note: this is the foundational ``why $\sigma$ exists'' answer for the
% course. Distractors (A), (B), (C), (E) each name something true {\em about} networks
% with $\sigma$, but they describe consequences rather than the structural cause. The
% test of whether a student has understood the role of $\sigma$ is whether they can
% articulate, without prompting, that affine-composed-with-affine is affine and that
% this is what stacking-without-$\sigma$ does.]
%
% \textbf{Trick answer:} B (D4 surface match: imports a backpropagation concern.
%   Affine maps are perfectly differentiable, and backpropagation through a network
%   with no activations works mechanically --- it just computes the gradient of an
%   affine map, which exists and is well-defined. The training pathology if you
%   removed $\sigma$ is not that gradients vanish but that the only function the
%   network {\em can} represent is affine. The student who picks (A) is reaching
%   for ``$\sigma$ is needed for training'' without checking whether removing $\sigma$
%   actually breaks training mechanics.)
% \textbf{Trick answer:} C (D6 shape-implies-property: confuses the role of nonlinearity
%   in general with the specific behavior of sigmoid. Sigmoid bounds outputs to
%   $(0,1)$; tanh bounds to $(-1,1)$; ReLU does not bound at all. Boundedness is a
%   feature of {\em particular} activations, not the structural reason any nonlinearity
%   is needed. A network using ReLU still escapes the affine-collapse problem despite
%   its activations being unbounded above.)
% \textbf{Trick answer:} A (D7 tautology: ``UAT does not apply'' restates the question
%   in different vocabulary. The Universal Approximation Theorem {\em assumes} a
%   nonlinear activation as a hypothesis, so saying ``without $\sigma$ UAT fails''
%   is just naming UAT's hypothesis, not explaining why $\sigma$ is needed. The
%   genuine answer addresses the structural cause that makes UAT's hypothesis
%   necessary in the first place.)
% \textbf{Trick answer:} D (D5 naive-prior: locates the problem in the world rather
%   than in the network. Even if every target function in the universe were linear,
%   removing $\sigma$ would still mean the network's depth and width are wasted ---
%   you would have a 50-layer linear regressor that does the same thing as a
%   1-layer linear regressor. The issue is intrinsic to the network architecture,
%   not contingent on the target.)
\divider

{\bf PROBLEM 28:}
% WEEK = 8
% TOPICS = Complex eigenvalues, oscillatory ODE solutions, Euler's formula
% DIFFICULTY = Core
%
A real matrix $A$ has a complex conjugate pair of eigenvalues $\lambda = \alpha \pm i\beta$ with $\beta \neq 0$. The system $d\mathbf{x}/dt = A\mathbf{x}$ has real basis solutions involving $e^{\alpha t} \cos(\beta t)$ and $e^{\alpha t} \sin(\beta t)$. Which of the following is the most fundamental reason these trigonometric functions appear?

\begin{enumerate}[(A)]
\item Euler's formula $e^{i\beta t} = \cos(\beta t) + i\sin(\beta t)$ converts the complex exponential solution into trigonometric form.
\item Because $A$ is real, complex eigenvalues must occur in conjugate pairs, forcing the imaginary parts of solutions to cancel.
\item The $2 \times 2$ real block associated with $\alpha \pm i\beta$ has the form $\begin{bmatrix} \alpha & -\beta \\ \beta & \alpha \end{bmatrix}$, which is a scaled rotation matrix.
\item Complex numbers do not appear in physically meaningful solutions, so they must be replaced by trigonometric functions.
\item The matrix exponential's Taylor series collapses into the series for sine and cosine when applied to the Jordan block.
\end{enumerate}

% \textbf{Correct Answer:} A
%
% \textbf{Explanation:} The complex-valued function $e^{\lambda t}\mathbf{v} = e^{(\alpha + i\beta)t}\mathbf{v}$ is a genuine solution to $d\mathbf{x}/dt = A\mathbf{x}$. Applying Euler's formula, $e^{(\alpha + i\beta)t} = e^{\alpha t}(\cos(\beta t) + i\sin(\beta t))$, so the complex solution splits into a real part and an imaginary part, each of which is itself a real solution by linearity. The trigonometric functions appear precisely because Euler's formula \emph{is} the bridge between complex exponentials and oscillation. Every other appearance of $\cos(\beta t)$ and $\sin(\beta t)$ in this setting — in the $2\times 2$ block formula, in the matrix-exponential Taylor series, in the structure of damped oscillators — derives from this single identity.
%
% The conjugate-pair structure (B) explains why we end up with \emph{real} solutions, but does not by itself explain why those solutions are \emph{trigonometric}. The block-form (C) and Taylor-series (D) facts are correct \emph{consequences}: they are how Euler's formula manifests in matrix language, not the underlying reason. Pedagogically, when students see ``where does $\cos(\beta t)$ come from?'', the answer should be Euler's formula first; the matrix-exponential derivation is a verification of what Euler already guarantees.
%
% \textbf{Trick answer:} B (correct fact, wrong question — conjugate pairing explains real-valuedness, not trigonometric structure. A classic D4 surface-match: ``complex conjugate'' is the right vocabulary for a different question.)
% \textbf{Trick answer:} C (true fact about the real Jordan block, but the rotation interpretation is itself a consequence of Euler's formula applied to $e^{Jt}$. Catches students who pattern-match on ``rotation'' without asking where the rotation came from.)
% \textbf{Partial credit:} E (the Taylor-series derivation is a correct verification, but it appeals to the structural identity $N^2 = -\beta^2 I$ which is itself the matrix shadow of $i^2 = -1$ — i.e., it reduces back to Euler's formula at the scalar level. A sophisticated student might pick this; partial credit is reasonable.)
% \textbf{Trick answer:} D (the naive-prior D5: confuses physical motivation with mathematical justification. Complex numbers do appear in solutions; we extract real solutions by linearity, not by fiat.)
%
% \textbf{Partial credit:} E (captures a correct derivation pathway, but reduces to Euler's formula one level deeper)

\divider

{\bf PROBLEM 29:}
% WEEK = 5
% TOPICS = Gram-Schmidt, partial-span invariance, orthogonal basis construction
%
Let $\{\mathbf{v}_1, \mathbf{v}_2, \ldots, \mathbf{v}_n\}$ be a linearly
independent set in an inner product space $V$, and let $\{\mathbf{u}_1,
\mathbf{u}_2, \ldots, \mathbf{u}_n\}$ be the orthogonal set produced by
applying the Gram-Schmidt process (in that order).  Which of the
following must be true for every $k$ with $1 \leq k \leq n$?
%
\begin{enumerate}[(A)]
\item $\mathbf{u}_k = \mathbf{v}_k$
\item $\norm{\mathbf{u}_k} = \norm{\mathbf{v}_k}$
\item $\mathrm{span}\{\mathbf{u}_1, \ldots, \mathbf{u}_k\} = \mathrm{span}\{\mathbf{v}_1, \ldots, \mathbf{v}_k\}$
\item $\mathbf{u}_k$ is a positive scalar multiple of $\mathbf{v}_k$
\item $\mathbf{u}_k \perp \mathbf{v}_k$
\end{enumerate}
%
% \textbf{Correct Answer:} C
%
% \textbf{Explanation:} The defining feature of Gram-Schmidt is that
% it builds the orthogonal vectors recursively by subtracting from
% $\mathbf{v}_k$ its projections onto $\mathbf{u}_1, \ldots,
% \mathbf{u}_{k-1}$:
% \[
%   \mathbf{u}_k = \mathbf{v}_k - \sum_{j=1}^{k-1} \frac{\langle \mathbf{v}_k, \mathbf{u}_j\rangle}{\langle \mathbf{u}_j, \mathbf{u}_j\rangle}\, \mathbf{u}_j.
% \]
% Each $\mathbf{u}_k$ lies in $\mathrm{span}\{\mathbf{v}_1, \ldots,
% \mathbf{v}_k\}$ (it is $\mathbf{v}_k$ plus a combination of previous
% $\mathbf{u}_j$'s, each of which is itself in that span), and
% conversely each $\mathbf{v}_k$ can be recovered from the
% $\mathbf{u}$'s.  So the partial spans match for every $k$.  This is
% why Gram-Schmidt produces a {\em compatible} orthogonal basis — it
% orthogonalizes without disturbing the nested-subspace structure of
% the original list.
%
% [Geometric picture: at step $k$, Gram-Schmidt takes $\mathbf{v}_k$
% and removes its component inside the subspace already spanned by
% $\mathbf{v}_1, \ldots, \mathbf{v}_{k-1}$.  The leftover is
% $\mathbf{u}_k$, which lives in the same partial span as
% $\mathbf{v}_k$ but is now orthogonal to everything that came before.
% This nested-span invariance is exactly the structure that makes the
% QR decomposition $A = QR$ work, with $R$ upper-triangular precisely
% because $\mathbf{v}_k$ involves only $\mathbf{u}_1, \ldots,
% \mathbf{u}_k$.]
%
% \textbf{Trick answer:} A (only true at $k=1$; in general
%   $\mathbf{u}_k \neq \mathbf{v}_k$ because Gram-Schmidt subtracts
%   the projections onto previous vectors.  Naive-prior answer for a
%   student who treats Gram-Schmidt as ``relabeling.'')
% \textbf{Trick answer:} B (norms are not preserved; subtracting
%   projections strictly shortens the vector unless $\mathbf{v}_k$ is
%   already orthogonal to the previous span.  Surface match on the
%   QR fact that the diagonal of $R$ encodes these length changes.)
% \textbf{Trick answer:} E (mixes up which vectors are orthogonal to
%   which.  Gram-Schmidt makes $\mathbf{u}_k \perp \mathbf{u}_j$ for
%   $j < k$, not $\mathbf{u}_k \perp \mathbf{v}_k$.  In fact
%   $\langle \mathbf{u}_k, \mathbf{v}_k \rangle = \norm{\mathbf{u}_k}^2
%   > 0$ since $\mathbf{v}_k = \mathbf{u}_k + (\text{stuff orthogonal
%   to } \mathbf{u}_k)$.)
% \textbf{Trick answer:} D (true at $k=1$ only; ignores that for $k
%   \geq 2$, $\mathbf{u}_k$ is $\mathbf{v}_k$ {\em minus} other
%   vectors, not a scalar multiple.  Twin pair to (A) — both stop
%   the recursion at the first step.)
%
\divider

{\bf PROBLEM 30:}
% WEEK = 3
% TOPICS = FTLA, coimage, image, cokernel, quotient spaces, dimension count
% DIFFICULTY = Deeper
%
Let $T: V \to W$ be a rank 3 linear transformation between finite-dimensional vector 
spaces, with $\dim(V) = 7$, $\dim(W) = 4$. 
Which of the following must be TRUE?
\begin{enumerate}[(A)]
\item $\dim(\ker T) = 3$ and $\mathrm{coim}(T) \cong \mathrm{im}(T)$
\item $\dim(\ker T) = 4$ and $\mathrm{coim}(T) \cong \ker(T)$
\item $\dim(\mathrm{coim}\, T) = 4$ and $\mathrm{coker}(T) \cong \mathbb{R}^1$
\item $\dim(\mathrm{coim}\, T) = 3$ and $\mathrm{coker}(T) \cong \mathbb{R}^1$
\item $\dim(\ker T) = 4$ and $\mathrm{coker}(T) \cong \ker(T)$
\end{enumerate}
% \textbf{Correct Answer:} D
%
% \textbf{Explanation:} Three facts are in play:
% (i) Rank-nullity: $\dim(\ker T) = \dim V - \dim(\mathrm{im}\, T) = 7 - 3 = 4$.
% (ii) FTLA punchline: $\mathrm{coim}(T) \cong \mathrm{im}(T)$, so 
%      $\dim(\mathrm{coim}\, T) = 3$.
% (iii) Cokernel dimension: $\dim(\mathrm{coker}\, T) = \dim W - \dim(\mathrm{im}\, T) 
%      = 4 - 3 = 1$, so $\mathrm{coker}(T) \cong \mathbb{R}^1$.
% Option (D) is the only choice consistent with all three.
%
% [Teaching remark: the FTLA gives a natural isomorphism 
% $\mathrm{coim}(T) \cong \mathrm{im}(T)$, but \emph{not} 
% $\mathrm{coker}(T) \cong \ker(T)$ in general — those have the same 
% dimension only when $\dim V = \dim W$, and even then the isomorphism 
% is not natural. This distinction is the FTLA spine.]
%
% \textbf{Partial credit:} C ($\mathrm{coker}(T) \cong \mathbb{R}^1$ correct, 
%   but $\dim(\mathrm{coim}\, T) = 4$ confuses $\mathrm{coim}\, T = V/\ker T$ 
%   with the wrong dimension count — the student remembered the 
%   quotient construction but lost the FTLA equivalence 
%   $\mathrm{coim}(T) \cong \mathrm{im}(T)$. Twin pair to (D).)
% \textbf{Trick answer:} B ($\dim \ker T = 4$ correct, but pairs it with 
%   the wrong FTLA isomorphism $\mathrm{coim}(T) \cong \ker(T)$. 
%   Coim is naturally isomorphic to \emph{image}, not kernel — this 
%   is the FTLA confusion that swaps the two non-trivial pieces of 
%   the four-subspace picture.)
% \textbf{Trick answer:} A (Has the correct FTLA isomorphism 
%   $\mathrm{coim}(T) \cong \mathrm{im}(T)$, but pairs it with the wrong 
%   kernel dimension $\dim \ker T = 3$ — confusing $\dim \ker T$ with 
%   $\mathrm{rank}(T)$, or equivalently using $\dim W - \mathrm{rank}$ 
%   instead of $\dim V - \mathrm{rank}$.)
% \textbf{Trick answer:} E ($\dim \ker T = 4$ correct, but pairs it with 
%   the false isomorphism $\mathrm{coker}(T) \cong \ker(T)$. This is the 
%   classic FTLA-misremembering trap: dimension counts can coincide for 
%   square-domain/codomain cases, but here $\dim \ker T = 4 \neq 1 = 
%   \dim \mathrm{coker}\, T$.)
\divider

{\bf PROBLEM 31:}
% WEEK = 7, 8 SYNTHESIS
% TOPICS = Diagonalizability, distinct eigenvalues, sufficient vs. necessary conditions, eigenspace dimension
%
Let $A$ be an $n \times n$ matrix with real entries. Which of the following statements about $A$ is most true?

\begin{enumerate}[(A)]
\item $A$ is diagonalizable if and only if every eigenvalue of $A$ is real.
\item If $A$ is diagonalizable, then $A$ has $n$ distinct eigenvalues, but the converse does not hold.
\item If $A$ has $n$ distinct real eigenvalues, then $A$ is diagonalizable, but the converse does not hold.
\item $A$ is diagonalizable if and only if $A$ has $n$ distinct real eigenvalues.
\item $A$ is diagonalizable if and only if $\det(A) \neq 0$.
\end{enumerate}

% \textbf{Correct Answer:} C
%
% \textbf{Explanation:} Having $n$ distinct eigenvalues is \emph{sufficient} for diagonalizability: distinct eigenvalues guarantee linearly independent eigenvectors, giving a full eigenbasis. But it is \emph{not necessary}: the identity $I_n$ has the single repeated eigenvalue $1$ (multiplicity $n$) and is already diagonal. More generally, $A$ is diagonalizable iff for every eigenvalue, geometric multiplicity equals algebraic multiplicity --- i.e., the eigenspaces together span $\mathbb{R}^n$ (or $\mathbb{C}^n$). Distinct eigenvalues force this automatically; repeated eigenvalues may or may not.
%
% [Geometric picture: distinct eigenvalues $\Rightarrow$ each eigenspace is $1$-dim and they pile up to $n$ independent directions. Repeated eigenvalues only obstruct diagonalization when the corresponding eigenspace is too small --- a Jordan block. The implication runs one way only.]
%
% \textbf{Trick answer:} D (asserts the converse, the classic confusion of sufficient with necessary; the canonical counterexample is $I_n$ or any $cI$.)
% \textbf{Trick answer:} B (reverses the implication; this would mean $I$ is non-diagonalizable, which is absurd since $I$ is already diagonal.)
% \textbf{Trick answer:} A (reality of eigenvalues is independent of diagonalizability: a $2\times 2$ rotation by $\theta \notin \pi\mathbb{Z}$ has complex eigenvalues and is diagonalizable over $\mathbb{C}$; a defective Jordan block can have all real eigenvalues and not be diagonalizable.)
% \textbf{Trick answer:} E (confuses diagonalizability with invertibility; $I$ is diagonalizable and invertible, the zero matrix is diagonalizable but singular, and a defective $2\times 2$ Jordan block at $\lambda = 1$ is invertible but not diagonalizable.)

\divider

{\bf PROBLEM 32:}
% WEEK = 10
% TOPICS = Polar decomposition, geometric meaning of A = QH, relationship to SVD
% DIFFICULTY = Deeper
A square matrix $A \in \mathbb{R}^{n \times n}$ admits a polar factorization $A = QH$, where $Q$ is orthogonal and $H$ is symmetric positive semidefinite.

Which statement best describes the geometric content of this decomposition?

\begin{enumerate}[(A)]
\item $Q$ rotates $\mathbb{R}^n$ into the principal-axis frame, and $H$ stretches each coordinate by the corresponding singular value.
\item $H = A^T A$ and $Q = A H^{-1}$, with $H$ symmetric positive semidefinite and $Q$ orthogonal.
\item $H$ stretches by the eigenvalues of $A$ along the eigenvectors of $A$, and $Q$ realigns the result.
\item $H$ is the symmetric part $(A + A^T)/2$ and $Q$ is the orthogonal part $(A - A^T)/2$ of $A$.
\item $H$ stretches along $n$ mutually orthogonal directions by the singular values of $A$, and $Q$ then rotates the stretched result.
\end{enumerate}

% \textbf{Correct Answer:} E
%
% \textbf{Explanation:} Polar decomposition writes the action of $A$ as
% \emph{stretch first, rotate second}. The symmetric PSD factor
% $H = \sqrt{A^T A}$ is orthogonally diagonalizable; its eigenvalues
% are exactly the singular values $\sigma_i$ of $A$, and its
% eigenvectors are the right singular vectors $\mathbf{v}_i$. So $H$
% acts on $\mathbb{R}^n$ by stretching along $n$ mutually orthogonal
% axes (the $\mathbf{v}_i$) by factors $\sigma_i$. Then $Q$ rigidly
% rotates the result into the codomain. This is the
% \textbf{fundamental} description: $A$ is a stretch composed with a
% rotation, with singular values controlling the stretch.
%
% Cross-check via SVD: $A = U \Sigma V^T = (UV^T)(V \Sigma V^T) = QH$,
% with $Q = UV^T$ and $H = V \Sigma V^T$.
%
% \textbf{Partial credit:} A (right structural intuition --- rotate
% then stretch --- but reverses the order; this is the \emph{left}
% polar decomposition $A = H'Q$, a different factorization with a
% different $H'$). One rung below correct.
%
% \textbf{Trick answer:} C (D4 surface match: confuses singular values
% with eigenvalues and right singular vectors with eigenvectors;
% $H$'s eigenvalues are $\sigma_i(A)$, not $\lambda_i(A)$, and
% these agree only when $A$ is symmetric PSD).
%
% \textbf{Trick answer:} B (D7 convention/notation collision: gets
% $H = \sqrt{A^T A}$ wrong as $H = A^T A$. The matrix $A^T A$ has
% eigenvalues $\sigma_i^2$, not $\sigma_i$, so this $H$ stretches by
% the {\em squares} of the singular values --- the wrong square root.
% Equivalently, $A H^{-1} = A (A^T A)^{-1}$ is not orthogonal in
% general; the formula for $Q$ inherits the same error).
%
% \textbf{Trick answer:} D (D5 naive-prior: imports the
% symmetric/skew-symmetric additive decomposition, which has nothing
% to do with polar form; the polar decomposition is multiplicative,
% not additive, and $(A - A^T)/2$ is skew-symmetric, not orthogonal).

\divider

{\bf PROBLEM 33:}
% WEEK = 3
% TOPICS = Linear transformations, isomorphism, injectivity, surjectivity, cokernel, FTLA
% DIFFICULTY = Deeper
%
Let $T: V \to W$ be a linear transformation between finite-dimensional 
vector spaces with $\dim(V) = \dim(W)$. Which of the following statements 
must be TRUE?

\begin{enumerate}[(A)]
\item The cokernel of $T$ is a subspace of $W$
\item If $T$ is injective, then $T$ is an isomorphism
\item If $\ker(T) = \{\mathbf{0}\}$, then $\mathrm{im}(T) = W$
\item Both (B) and (C)
\item All of the above
\end{enumerate}

% \textbf{Correct Answer:} D
%
% \textbf{Explanation:} Under the hypothesis $\dim(V) = \dim(W)$, injectivity 
% and surjectivity are equivalent — and either, alone, implies $T$ is an 
% isomorphism. By rank-nullity, $\dim(\mathrm{im}\, T) = \dim(V) - 
% \dim(\ker T)$. If $T$ is injective then $\dim(\ker T) = 0$, so 
% $\dim(\mathrm{im}\, T) = \dim(V) = \dim(W)$, forcing $\mathrm{im}(T) = W$ 
% and giving surjectivity (and thus an isomorphism). Statement (C) says 
% exactly this — injective implies surjective — using the kernel 
% characterization of injectivity. So (B) and (C) are equivalent statements 
% of the same fact, and the genuine answer is (D).
%
% [Teaching remark: the equal-dimension hypothesis collapses injectivity 
% and surjectivity into a single condition, just as it does for square 
% matrices (where invertibility, injectivity, and surjectivity all 
% coincide). Without equal dimensions, all three properties can come apart 
% — the differentiation operator $\mathcal{D}: \mathcal{P}_3 \to \mathcal{P}_2$ 
% from Q1 P5 is surjective but not injective. Recognizing that (B) and (C) 
% describe the same fact via the FTLA is the conceptual payoff.]
%
% \textbf{Trick answer:} A (categorical confusion: $\mathrm{coker}(T) = 
%   W / \mathrm{im}(T)$ is a \emph{quotient} of $W$, not a subspace. The 
%   trap is the surface intuition ``cokernel lives on the $W$-side, so it 
%   must sit inside $W$.'' This conflates the quotient construction with 
%   subspace inclusion — a Wk-3 conceptual error worth catching. Note that 
%   $\mathrm{coker}(T)$ \emph{is} isomorphic to a subspace of $W$ when an 
%   inner product is available, but as defined it is a quotient, not a 
%   subspace.)
% \textbf{Partial credit:} B (recognizes the injective $\Rightarrow$ 
%   isomorphism direction — captures one ladder rung but misses that (C) 
%   is the equivalent kernel-based statement of the same fact. One rung 
%   below correct.)
% \textbf{Partial credit:} C (recognizes that trivial kernel forces full 
%   image under the dimension hypothesis — captures the FTLA dimension 
%   count but misses that (B) is the equivalent statement in 
%   injectivity language. Twin pair to (B); same ladder rung.)
% \textbf{Trick answer:} E (requires knowing that (A) is false.)

\divider



{\bf PROBLEM 34:}
% WEEK = 5, 6, 12 SYNTHESIS
% TOPICS = Perceptron geometry, decision boundary, normal vector,
%          signed distance, projection, dot product as geometric quantity
% DIFFICULTY = Core (conceptual, synthesis)
A perceptron with weights $\mathbf{w} \in \mathbb{R}^n$, bias $b \in \mathbb{R}$, and sigmoid activation $\sigma$ computes $y = \sigma(\mathbf{w}^T \mathbf{x} + b)$, with decision boundary $\{\mathbf{x} : \mathbf{w}^T \mathbf{x} + b = 0\}$.
%
Which statement best and most fundamentally describes the geometric role of $\mathbf{w}$?

\begin{enumerate}[(A)]
\item $\mathbf{w}$ points from the origin toward the centroid of the training inputs whose target label is close to $1$.
\item $\mathbf{w}$ lies in the decision boundary, since $\mathbf{w}$ and $\mathbf{x}$ are vectors in the same space $\mathbb{R}^n$.
\item $\mathbf{w}$ is normal to the decision boundary, and $\mathbf{w}^T \mathbf{x} + b$ equals $\norm{\mathbf{w}}$ times the signed distance from $\mathbf{x}$ to that boundary.
\item $\mathbf{w}$ determines the decision boundary.
\item $\mathbf{w}$ is the direction of steepest ascent of the pre-activation $\mathbf{x} \mapsto \mathbf{w}^T \mathbf{x} + b$.
\end{enumerate}

% \textbf{Correct Answer:} C
%
% \textbf{Explanation:} Two facts together pin down the geometric role of $\mathbf{w}$.
%
% (i) The decision boundary is the level set $\{\mathbf{x} : \mathbf{w}^T \mathbf{x} + b = 0\}$,
% an affine hyperplane. For any two points $\mathbf{x}_1, \mathbf{x}_2$ on the boundary,
% $\mathbf{w}^T(\mathbf{x}_1 - \mathbf{x}_2) = 0$, so $\mathbf{w}$ is orthogonal to every direction
% along the boundary --- i.e., $\mathbf{w}$ is a \emph{normal vector} to the hyperplane.
%
% (ii) Decompose any input as $\mathbf{x} = \mathbf{x}_\parallel + t \cdot \mathbf{w} / \norm{\mathbf{w}}$,
% where $\mathbf{x}_\parallel$ is the projection of $\mathbf{x}$ onto the boundary and $t$ is
% its signed distance to the boundary along $\mathbf{w}$. Then
% $\mathbf{w}^T \mathbf{x} + b = \mathbf{w}^T \mathbf{x}_\parallel + b + t \norm{\mathbf{w}} = t \norm{\mathbf{w}}$,
% since $\mathbf{w}^T \mathbf{x}_\parallel + b = 0$ on the boundary. So the pre-activation is
% literally $\norm{\mathbf{w}}$ times the signed distance from $\mathbf{x}$ to the decision boundary.
%
% This is the synthesis: the perceptron is a geometric construction in disguise. The dot
% product $\mathbf{w}^T \mathbf{x}$ from Wk 5, the orthogonal-projection picture from Wk 6,
% and the affine-hyperplane / kernel picture from Wk 1--3 all collapse together. ``Classify
% $\mathbf{x}$'' means ``compute the signed distance from $\mathbf{x}$ to a hyperplane and
% squash through $\sigma$.''
%
% \textbf{Trick answer:} D but misses that the pre-activation \emph{magnitude} also has metric
%   meaning, namely $\norm{\mathbf{w}}$ times signed distance. One rung below (B) on the ladder.
% \textbf{Trick answer:} E (true, but a \emph{calculus} consequence of the linear-algebraic fact
%   in (B): once you know the level sets of $\mathbf{w}^T \mathbf{x} + b$ are hyperplanes normal
%   to $\mathbf{w}$, ``steepest ascent direction is $\mathbf{w}$'' follows. The gradient fact
%   describes a \emph{consequence} of $\mathbf{w}$'s role; it does not capture what $\mathbf{w}$
%   geometrically \emph{is}.)
% \textbf{Trick answer:} B (input/weight-space confusion. Living in the same ambient $\mathbb{R}^n$
%   does not put $\mathbf{w}$ on the boundary --- $\mathbf{w}$ generically does not satisfy
%   $\mathbf{w}^T \mathbf{w} + b = 0$. The fact that $\mathbf{w}$ is normal to the boundary is
%   precisely the opposite of lying in it.)
% \textbf{Trick answer:} A (folk picture: the weight vector ``points toward the positive class.''
%   Directionally suggestive but mathematically false in general --- the centroid of positive
%   examples need not lie along $\mathbf{w}$, and after training $\mathbf{w}$ is determined by
%   margin geometry, not class centroids. This story is a half-remembered description of LDA
%   or a class-mean classifier, not a perceptron.)

\divider

{\bf PROBLEM 35:}
% WEEK = 6
% TOPICS = Orthogonal projection, best approximation, closest point in a subspace
%
Let $V$ be a finite-dimensional inner product space and let $W \subseteq V$
be a subspace.  For $\mathbf{v} \in V$, let $\Pi_W(\mathbf{v})$ denote
the orthogonal projection of $\mathbf{v}$ onto $W$.  Which of the
following characterizes $\Pi_W(\mathbf{v})$?
%
\begin{enumerate}[(A)]
\item The unique vector $\mathbf{w} \in W$ with $\norm{\mathbf{w}} = \norm{\mathbf{v}}$
\item The unique vector $\mathbf{w} \in W$ minimizing $\norm{\mathbf{w}}$
\item The unique vector $\mathbf{w} \in V$ such that $\mathbf{v} - \mathbf{w} \in W^\perp$
\item The unique vector $\mathbf{w} \in W$ minimizing $\norm{\mathbf{v} - \mathbf{w}}$
\item The unique vector $\mathbf{w} \in W$ such that $\mathbf{v} - \mathbf{w} \in W$
\end{enumerate}
%
% \textbf{Correct Answer:} D
%
% \textbf{Explanation:} The orthogonal projection $\Pi_W(\mathbf{v})$ is
% the {\em closest point in $W$ to $\mathbf{v}$}: it is the unique
% $\mathbf{w} \in W$ minimizing $\norm{\mathbf{v} - \mathbf{w}}$.
% Equivalently, it is the unique $\mathbf{w} \in W$ such that the
% residual $\mathbf{v} - \mathbf{w}$ lies in $W^\perp$.
%
% [Geometric picture: drop a perpendicular from $\mathbf{v}$ to $W$; the
% foot of the perpendicular is $\Pi_W(\mathbf{v})$.  The residual
% $\mathbf{v} - \Pi_W(\mathbf{v})$ is the perpendicular itself, lying
% in $W^\perp$.]
%
% \textbf{Trick answer:} B (the smallest-norm vector in $W$ is
%   $\mathbf{0}$ for any subspace; this characterization ignores
%   $\mathbf{v}$ entirely.  Naive-prior surface match on "minimization.")
% \textbf{Trick answer:} E (asks for $\mathbf{v} - \mathbf{w} \in W$
%   instead of $W^\perp$ — the orthogonality condition flipped to the
%   wrong subspace.  Twin pair to (D) at the subspace level.)
% \textbf{Trick answer:} C (correct orthogonality condition, but
%   $\mathbf{w}$ is allowed to range over $V$ rather than constrained
%   to $W$.  Without the constraint $\mathbf{w} \in W$ the equation
%   $\mathbf{v} - \mathbf{w} \in W^\perp$ has many solutions — every
%   element of $\Pi_W(\mathbf{v}) + W^\perp$ — so uniqueness fails.
%   Twin pair to (C); each loses one of the two defining conditions.)
% \textbf{Trick answer:} A (norm-preservation: a property of
%   {\em isometries}, not projections.  Projections shorten vectors
%   in general — surface match on "$\mathbf{w}$ relates to $\mathbf{v}$
%   by some norm equation.")
%
\divider

{\bf PROBLEM 36:}
% WEEK = 7
% TOPICS = Linear ODE systems, matrix exponential, eigencomponent solution, initial conditions
%
Let $A$ be a $2 \times 2$ matrix with eigenvalues $\lambda_1 = 2$ and $\lambda_2 = -1$ and corresponding eigenvectors
$\mathbf{v}_1 = \begin{pmatrix} 1 \\ 1 \end{pmatrix},  \mathbf{v}_2 = \begin{pmatrix} 1 \\ -1 \end{pmatrix}$.
What is the solution $\mathbf{x}(t)$ to the initial value problem
\[
\frac{d\mathbf{x}}{dt} = A\mathbf{x}, \qquad \mathbf{x}(0) = \begin{pmatrix} 3 \\ 1 \end{pmatrix}.
\]


\begin{enumerate}[(A)]
\item $2 e^{2t} \begin{pmatrix} 1 \\ 1 \end{pmatrix} + e^{-t} \begin{pmatrix} 1 \\ -1 \end{pmatrix}$
\item $2 e^{2t} \begin{pmatrix} 1 \\ 1 \end{pmatrix} + e^{t} \begin{pmatrix} 1 \\ -1 \end{pmatrix}$
\item $e^{2t} \begin{pmatrix} 3 \\ 1 \end{pmatrix} + e^{-t} \begin{pmatrix} 3 \\ 1 \end{pmatrix}$
\item $e^{2t} \begin{pmatrix} 1 \\ 1 \end{pmatrix} + 2 e^{-t} \begin{pmatrix} 1 \\ -1 \end{pmatrix}$
\item $3 e^{2t} \begin{pmatrix} 1 \\ 1 \end{pmatrix} + e^{-t} \begin{pmatrix} 1 \\ -1 \end{pmatrix}$
\end{enumerate}

% \textbf{Correct Answer:} A
%
% \textbf{Explanation:} For diagonalizable $A$, the general solution to $\dot{\mathbf{x}} = A\mathbf{x}$ is
% \[
% \mathbf{x}(t) = c_1 e^{\lambda_1 t}\mathbf{v}_1 + c_2 e^{\lambda_2 t}\mathbf{v}_2,
% \]
% where the constants $c_1, c_2$ are the coordinates of $\mathbf{x}_0$ in the eigenbasis. Solving $\mathbf{x}(0) = c_1 \mathbf{v}_1 + c_2 \mathbf{v}_2 = \begin{pmatrix} 3 \\ 1 \end{pmatrix}$:
% \[
% c_1 + c_2 = 3, \qquad c_1 - c_2 = 1 \quad \Longrightarrow \quad c_1 = 2, \; c_2 = 1.
% \]
% Hence $\mathbf{x}(t) = 2 e^{2t}\mathbf{v}_1 + e^{-t}\mathbf{v}_2$.
%
% [Equivalently: $\mathbf{x}(t) = e^{At}\mathbf{x}_0 = V e^{\Lambda t} V^{-1}\mathbf{x}_0$. The vector $V^{-1}\mathbf{x}_0 = (c_1, c_2)^T$ is the initial condition expressed in the eigenbasis; $e^{\Lambda t}$ scales each coordinate by $e^{\lambda_i t}$; $V$ converts back to standard coordinates.]
%
% \textbf{Partial credit:} D (correct solution form and correct eigenvalue placement, but swaps the eigencomponent coefficients --- gets $c_1 = 1, c_2 = 2$ instead of $c_1 = 2, c_2 = 1$. Right framework, arithmetic slip in solving the $2 \times 2$ coordinate system.)
% \textbf{Trick answer:} E (uses the entries of $\mathbf{x}_0 = (3,1)^T$ directly as the eigencomponent coefficients $c_1 = 3, c_2 = 1$, skipping the eigenbasis change of coordinates. This is the most common mistake: forgetting that $c_i$ are coordinates \emph{in the eigenbasis}, not coordinates of $\mathbf{x}_0$ in the standard basis.)
% \textbf{Trick answer:} B (sign error on the second eigenvalue: $e^{t}$ instead of $e^{-t}$. Coefficients are correct but $\lambda_2 = -1$, not $+1$.)
% \textbf{Trick answer:} C (wholesale conceptual error: treats $\mathbf{x}_0$ as if it were the eigenvector for both modes, replacing $\mathbf{v}_1$ and $\mathbf{v}_2$ with $\mathbf{x}_0$. Misses the fundamental role of eigenvectors as the directions of independent exponential evolution.)

\divider

{\bf PROBLEM 37:}
% WEEK = 9
% TOPICS = Positive stochastic matrices, Markov chains, stationary distribution,
%          eigenvalue 1, Perron--Frobenius, dominant eigenvector
%
Let $P$ be a positive column-stochastic $n \times n$ matrix governing a Markov chain
\[
    \mathbf{x}_{k+1}=P\mathbf{x}_k
\]
on probability vectors in $\mathbb{R}^n$. A \emph{stationary distribution}
for $P$ is a probability vector $\boldsymbol{\pi}$ satisfying
\[
    P\boldsymbol{\pi}=\boldsymbol{\pi}.
\]

Which of the following is the best Perron--Frobenius description of the
stationary distribution?

\begin{enumerate}[(A)]
\item A nonzero vector in $\ker(P-I)$.
\item The limit $\lim_{k\to\infty} P^k\mathbf{x}_0$ for any probability vector $\mathbf{x}_0$.
\item A probability vector $\boldsymbol{\pi}$ such that $\boldsymbol{\pi}^T P = \boldsymbol{\pi}^T$.
\item An eigenvector of $P$ whose entries sum to $1$.
\item The probability eigenvector of $P$ corresponding to the dominant eigenvalue $\lambda=1$.
\end{enumerate}

% \textbf{Correct Answer:} E
%
% \textbf{Explanation:} Since $P$ is positive and column-stochastic,
% Perron--Frobenius applies cleanly.  The eigenvalue $\lambda=1$ is the
% dominant eigenvalue, and its eigenspace contains a positive eigenvector.
% Normalizing that positive eigenvector so that its entries sum to $1$
% gives the stationary distribution.
%
% Choice (D) captures both essential features: the vector is fixed by the
% dynamics, $P\boldsymbol{\pi}=\boldsymbol{\pi}$, and it is a probability
% vector.  The word ``dominant'' identifies the Perron--Frobenius role of
% the eigenvalue $\lambda=1$.
%
% \textbf{Partial credit:} A (recognizes the fixed-point equation
%   $P\boldsymbol{\pi}=\boldsymbol{\pi}$, equivalently
%   $\boldsymbol{\pi}\in\ker(P-I)$, but misses the probability
%   normalization/nonnegativity requirement.)
%
% \textbf{Trick answer:} C (row-vector convention.  Under the
%   column-stochastic convention used here, stationary distributions are
%   right eigenvectors of $P$, not left eigenvectors.)
%
% \textbf{Partial credit:} B (true under the positivity hypothesis, but
%   describes the long-term limiting behavior rather than the
%   Perron--Frobenius eigenvector description of what the stationary
%   distribution is.)
%
% \textbf{Trick answer:} D (requires entries summing to $1$, but does not
%   specify the eigenvalue $\lambda=1$ and does not require nonnegative
%   entries.  So it need not describe a probability vector fixed by the
%   dynamics.)
%
\divider

{\bf PROBLEM 38:}
% WEEK = 4
% TOPICS = Coordinate map, isomorphism via basis choice, abstract vs concrete vector spaces
%
Let $V$ be an $n$-dimensional vector space and fix
an ordered basis $\mathcal{B}$.  Consider the coordinate map
$\Phi_{\mathcal{B}} : V \to \mathbb{R}^n$ defined by
$\Phi_{\mathcal{B}}(\mathbf{v}) = [\mathbf{v}]_{\mathcal{B}}$.  Which
of the following must be true?
%
\begin{enumerate}[(A)]
\item $\Phi_{\mathcal{B}}$ is an isomorphism only if $V$ is literally $\mathbb{R}^n$ (not merely isomorphic to it).
\item $\Phi_{\mathcal{B}}$ is a linear transformation but generally not invertible.
\item $\Phi_{\mathcal{B}}$ is an isomorphism, and any two ordered bases $\mathcal{B}, \mathcal{C}$ produce the same coordinate map.
\item $\Phi_{\mathcal{B}}$ is an isomorphism, and a different ordered basis $\mathcal{C}$ gives a different coordinate map $\Phi_{\mathcal{C}}$.
\item $\Phi_{\mathcal{B}}$ is an isomorphism only if $\mathcal{B}$ is orthonormal.
\end{enumerate}
%
% \textbf{Correct Answer:} D
%
% \textbf{Explanation:} The coordinate map $\Phi_{\mathcal{B}}$ is
% linear (coordinates respect addition and scalar multiplication),
% injective (only the zero vector has all-zero coordinates), and
% surjective (every $(c_1, \ldots, c_n) \in \mathbb{R}^n$ is the
% coordinates of $c_1 \mathbf{b}_1 + \cdots + c_n \mathbf{b}_n$).
% So $\Phi_{\mathcal{B}}$ is an isomorphism for {\em any} ordered basis
% $\mathcal{B}$ of any $n$-dimensional $V$.  But the isomorphism is
% {\em not canonical} --- it depends on the basis choice.  For any two
% distinct ordered bases $\mathcal{B} \neq \mathcal{C}$, the coordinate
% maps $\Phi_{\mathcal{B}}$ and $\Phi_{\mathcal{C}}$ are different
% functions: they disagree on at least the basis vectors themselves,
% and in general they differ by the change-of-basis matrix
% $\Phi_{\mathcal{C}} = P \circ \Phi_{\mathcal{B}}$.
%
% [Conceptual takeaway: every finite-dimensional vector space ``becomes''
% $\mathbb{R}^n$ after a basis choice.  This is what justifies treating
% abstract vector spaces $V$ via numerical computations on coordinates ---
% but the price is that {\em the isomorphism depends on the basis}.
% Coordinate-free properties survive any basis choice; coordinate-dependent
% statements (like the matrix entries of a transformation) do not.]
%
% \textbf{Trick answer:} B (recognizes linearity but mis-remembers
%   invertibility --- the naive-prior answer for a student who has not
%   internalized that a basis is, by definition, exactly the data
%   needed to make the coordinate map bijective.)
% \textbf{Trick answer:} C (twin pair to (B); collapses the
%   basis-dependence of the coordinate map.  Tempts students who
%   remember ``coordinate maps are isomorphisms'' but treat the
%   isomorphism as canonical --- the same conceptual error as treating
%   $V \cong V^*$ via an unspecified isomorphism.)
% \textbf{Trick answer:} A (the naive-prior answer for the student who
%   thinks $\mathbb{R}^n$ is special; misses that coordinate maps
%   {\em produce} the identification $V \cong \mathbb{R}^n$ rather
%   than requiring it.  The clarification ``literally $\mathbb{R}^n$''
%   sharpens the trap by ruling out the tautological reading.)
% \textbf{Trick answer:} E (imports an inner-product condition that is
%   irrelevant to the coordinate map.  Surface match on ``orthonormal
%   bases are nicer'' --- true in inner product spaces for computation,
%   but the coordinate map exists and is an isomorphism without any
%   inner product structure.)
\divider

{\bf PROBLEM 39:}
% WEEK = 4
% TOPICS = Similarity, change of basis, matrix representations of a linear transformation
% DIFFICULTY = Core
Square matrices $A$ and $B$ of the same size satisfy $B = P^{-1} A P$ for some invertible matrix $P$. Which statement most fundamentally describes the relationship between $A$ and $B$?

\begin{enumerate}[(A)]
\item $A$ and $B$ have the same determinant and trace.
\item $A$ and $B$ have the same image as subspaces of $\mathbb{R}^n$.
\item $A$ and $B$ have the same eigenvalues.
\item $A$ is obtained from $B$ by a change of coordinates in the codomain only.
\item $A$ and $B$ represent the same linear transformation expressed in different bases.
\end{enumerate}

% \textbf{Correct Answer:} E
% \textbf{Explanation:} The equation $B = P^{-1} A P$ is the change-of-basis
% formula for a linear transformation $T : V \to V$. If $A = [T]_{\mathcal{B}}$
% and $\mathcal{C}$ is another basis with change-of-basis matrix $P$ (so columns
% of $P$ express the $\mathcal{C}$-vectors in $\mathcal{B}$-coordinates), then
% $[T]_{\mathcal{C}} = P^{-1} A P = B$. The fundamental content of similarity
% is therefore: \textbf{same transformation, different bases}. Every other true
% statement about similar matrices --- equal eigenvalues, equal determinant,
% equal trace, equal rank --- is a \emph{consequence} of this fact, because
% these quantities are intrinsic to the transformation $T$ itself and do not
% depend on the basis chosen to represent it.
%
% This is the conceptual seed of diagonalization: an eigenbasis is a basis
% in which the matrix representation is diagonal, because in that basis
% $T$ acts as a pure stretch along each axis.
%
% \textbf{Trick answer:} C (recognizes the strongest invariant of similarity
% --- eigenvalues are preserved because they are intrinsic to the transformation
% --- but identifies a consequence rather than the meaning.)
% \textbf{Trick answer:} A (recognizes two correct invariants of similarity,
% but a strictly weaker observation than (A), since equal eigenvalues already
% imply equal determinant and trace.)
% \textbf{Trick answer:} D (the change of basis affects \emph{both} the input
% and output sides, since the same basis is used for the domain and codomain
% of $T : V \to V$; this is why $P$ appears on both sides of $A$, not on one
% side. A change of coordinates ``on one side only'' would correspond to a
% different operation, not similarity.)
% \textbf{Trick answer:} B (the classic similarity trap: image, kernel, and
% column/row spaces depend on the choice of basis and are \emph{not} preserved
% under similarity. For example, $I$ and $P^{-1} I P = I$ are similar with
% identical image, but $A$ and $P^{-1} A P$ for general $A$ have images that
% differ as concrete subspaces of $\mathbb{R}^n$.)

\divider

{\bf PROBLEM 40:}
% WEEK = 8
% TOPICS = Jordan canonical form, algebraic multiplicity, geometric multiplicity, block-size partition
%
A $5 \times 5$ matrix $A$ has a single eigenvalue $\lambda = 3$ with algebraic multiplicity $5$ and geometric multiplicity $2$. Which of the following statements is most true about the Jordan canonical form of $A$?

\begin{enumerate}[(A)]
\item The Jordan form has exactly $2$ Jordan blocks, but their sizes are not determined by the given information.
\item The Jordan form has exactly $3$ Jordan blocks.
\item The Jordan form is diagonal, since all eigenvalues equal $3$.
\item The Jordan form is uniquely determined to be one $4 \times 4$ Jordan block plus one $1 \times 1$ block.
\item None of the above.
\end{enumerate}

% \textbf{Correct Answer:} A
%
% \textbf{Explanation:} For a single eigenvalue, the geometric multiplicity equals the \emph{number} of Jordan blocks, while the algebraic multiplicity equals the total number of rows occupied by those blocks (the sum of block sizes). With algebraic multiplicity $5$ and geometric multiplicity $2$, the Jordan form has exactly $2$ blocks whose sizes sum to $5$. The possible partitions of $5$ into $2$ positive parts are $5 = 4 + 1$ and $5 = 3 + 2$, giving two distinct (non-similar) Jordan forms consistent with the data:
% \[
% J_4(3) \oplus J_1(3) \quad \text{or} \quad J_3(3) \oplus J_2(3).
% \]
% Distinguishing them requires more refined information --- e.g., the dimensions of $\ker((A - 3I)^k)$ for $k = 2, 3, \ldots$, which count generalized-eigenspace filtration steps. So algebraic and geometric multiplicities pin down the \emph{count} of blocks but not their \emph{sizes}.
%
% [Rule of thumb: geometric multiplicity = number of blocks; algebraic multiplicity = total size. The partition of block sizes is a \emph{finer} invariant than the spectrum-plus-multiplicities pair, and it is precisely what the Jordan form adds to the eigenvalue picture.]
%
% \textbf{Partial credit:} D (correctly identifies one of the two valid Jordan forms and the right block count, but asserts uniqueness; this is the ``picked the largest possible block'' instinct, which is one consistent answer but not the only one.)
% \textbf{Trick answer:} B (D6 shape-implies-property: confuses ``deficiency'' (alg.\ mult. minus geom.\ mult. $= 5 - 2 = 3$, which counts the number of \emph{generalized} eigenvectors needed beyond ordinary ones) with the number of blocks. The number of blocks is the geometric multiplicity, $2$ --- not $3$.)
% \textbf{Trick answer:} C (asserts diagonalizability from a repeated eigenvalue without checking geometric multiplicity. With geom.\ mult.\ $= 2 \neq 5 =$ alg.\ mult., $A$ is not diagonalizable, so $A \neq 3I_5$. This is the P2030-C diagonalizability-gap trap.)
% \textbf{Trick answer:} E 

\divider

{\bf PROBLEM 41:}
% WEEK = 7, 8 SYNTHESIS
% TOPICS = Matrix exponential of a Jordan block, nilpotent decomposition,
%          origin of the polynomial ladder t^k e^(lambda t) in repeated-root ODE solutions
% DIFFICULTY = Deeper
%
Let $J = \lambda I + N$ be a $k \times k$ Jordan block with eigenvalue
$\lambda$ on the diagonal and a single superdiagonal of $1$'s, so that $N$
is the $k \times k$ nilpotent matrix with $N^k = 0$. The matrix exponential
$e^{Jt}$ controls the solution to $d\mathbf{x}/dt = J\mathbf{x}$.
Which of the following is the most accurate description of $e^{Jt}$, and
why?

\begin{enumerate}[(A)]
\item $e^{Jt} = e^{\lambda t}\!\left(I + tN + \tfrac{t^2}{2!}N^2 + \cdots + \tfrac{t^{k-1}}{(k-1)!}N^{k-1}\right)$, because $\lambda I$ and $N$ commute and $N$ is nilpotent of index $k$.
\item $e^{Jt} = e^{\lambda I t} + e^{Nt}$, since the matrix exponential is additive over a sum of commuting matrices.
\item $e^{Jt}$ is undefined whenever $J$ is non-diagonalizable, since the matrix exponential requires an eigenbasis to compute.
\item $e^{Jt} = e^{\lambda t}\!\left(I + tN + \tfrac{t^2}{2!}N^2 + \cdots\right)$ as an infinite series, since the superdiagonal contributions never terminate.
\item $e^{Jt} = e^{\lambda t} I$, since the diagonal entry $\lambda$ governs the exponential growth rate.
\end{enumerate}

% \textbf{Correct Answer:} A
%
% \textbf{Explanation:} The decomposition $J = \lambda I + N$ exhibits $J$
% as a sum of two \emph{commuting} matrices --- $\lambda I$ commutes with
% everything --- so the matrix exponential factors as
% $$e^{Jt} \;=\; e^{(\lambda I + N) t} \;=\; e^{\lambda I t}\, e^{N t} \;=\; e^{\lambda t}\, e^{Nt}.$$
% The first factor is the scalar $e^{\lambda t}$ times $I$; the second factor
% is the matrix exponential of the nilpotent $N$, which by definition is the
% power series $e^{Nt} = \sum_{j \geq 0} \frac{t^j}{j!} N^j$. Because $N^k = 0$,
% this series terminates after $k$ terms:
% $$e^{Nt} \;=\; I + tN + \tfrac{t^2}{2!} N^2 + \cdots + \tfrac{t^{k-1}}{(k-1)!} N^{k-1}.$$
% Multiplying by $e^{\lambda t}$ gives (C).
%
% This is the structural origin of the polynomial-times-exponential ladder
% $\{e^{\lambda t}, t e^{\lambda t}, t^2 e^{\lambda t}/2!, \ldots\}$ that
% appears as the basis of solutions to a higher-order linear ODE with a
% repeated root (P21). The factor $e^{\lambda t}$ comes from $\lambda I$;
% the polynomial factors $1, t, t^2/2!, \ldots, t^{k-1}/(k-1)!$ are the
% entries that $N$ generates as it climbs the superdiagonal. The Jordan
% block is the matrix avatar of the repeated-root behavior --- Jordan
% structure and the polynomial ladder are the same phenomenon viewed
% from two angles.
%
% [Conceptual through-line: matrix exponentials of diagonalizable matrices
% (Wk 7) involve only $e^{\lambda t}$ scalings on each eigenvector. The
% \emph{new} content of Wk 8 is what happens when there are not enough
% eigenvectors --- the missing directions are filled in by the nilpotent
% $N$, and matrix exponentiation reaches into those directions through
% the polynomial $t^j N^j / j!$ ladder. (C) is the single formula that
% encodes both the eigenvalue scaling and the Jordan-chain ladder.]
%
% \textbf{Trick answer:} E (D5 naive-prior: takes the diagonal-case answer
%   $e^{\lambda t} I$ and applies it to the non-diagonalizable case. This
%   would be correct if $N = 0$, but the whole point of a Jordan block is
%   that $N \neq 0$; ignoring $N$ ignores the entire Jordan-chain
%   structure. A student who picks (A) has not registered that
%   non-diagonalizable matrices behave differently under exponentiation.)
% \textbf{Trick answer:} B (D7 algebraic-rule misuse: invokes a false
%   ``additivity'' rule. The matrix exponential is multiplicative over
%   commuting summands, $e^{A+B} = e^A e^B$ when $AB = BA$ --- never
%   additive. The correct decomposition does start from
%   $J = \lambda I + N$ with commutativity, but produces a \emph{product}
%   $e^{\lambda I t} e^{Nt}$, not a sum. The student who picks (B) has
%   the right ingredients but the wrong combinator.)
% \textbf{Partial credit:} D (gets the structural picture right ---
%   $e^{Jt} = e^{\lambda t} e^{Nt}$ with $e^{Nt}$ as a series in $N$ ---
%   but fails to register that $N^k = 0$ truncates the series. The
%   student understands the decomposition but has not exploited
%   nilpotency. One rung below (C) on the ladder: same framework,
%   missing the finite-series payoff.)
% \textbf{Trick answer:} C (D5 naive-prior: confuses ``computation by
%   diagonalization'' with the existence of $e^{Jt}$. The matrix
%   exponential is defined for \emph{any} square matrix via its power
%   series; diagonalization is a convenient shortcut when available, not
%   a precondition. The Jordan-block formula in (C) is precisely the
%   tool that handles the non-diagonalizable case.)

\divider

{\bf PROBLEM 42:}
% WEEK = 11
% TOPICS = Correlation as cosine similarity, centered data, geometric interpretation,
%          inner product structure on feature columns
% DIFFICULTY = Deeper
%
A data matrix $X \in \mathbb{R}^{n \times d}$ has $n$ observations of $d$ features
arranged in columns $\mathbf{x}_1, \ldots, \mathbf{x}_d \in \mathbb{R}^n$, each
centered to have sample mean zero. Let $\rho_{ij}$ denote the Pearson correlation
coefficient between features $i$ and $j$. Which of the following best describes
the geometric meaning of $\rho_{ij}$?

\begin{enumerate}[(A)]
\item $\rho_{ij}$ is the cosine of the angle between $\mathbf{x}_i$ and $\mathbf{x}_j$, viewed as vectors in $\mathbb{R}^n$.
\item $\rho_{ij}$ is the length of the projection of $\mathbf{x}_i$ onto $\mathbf{x}_j$.
\item $\rho_{ij}$ is the cosine of the angle between the $i$th and $j$th rows of $X$, viewed as vectors in $\mathbb{R}^d$.
\item $\rho_{ij}$ is the cosine of the angle between $\mathbf{x}_i$ and $\mathbf{x}_j$ only after both have been further rescaled to have unit norm.
\item $\rho_{ij}$ equals the inner product $\langle \mathbf{x}_i, \mathbf{x}_j \rangle$ scaled by the number of observations.
\end{enumerate}

% \textbf{Correct Answer:} A
%
% \textbf{Explanation:} For centered feature columns, the Pearson correlation
% coefficient is precisely the cosine similarity of the columns viewed as
% vectors in $\mathbb{R}^n$:
% $$\rho_{ij} = \frac{\sum_k x_{ki} x_{kj}}{\sqrt{\sum_k x_{ki}^2}\sqrt{\sum_k x_{kj}^2}}
%   = \frac{\langle \mathbf{x}_i, \mathbf{x}_j \rangle}{\norm{\mathbf{x}_i}\,\norm{\mathbf{x}_j}}
%   = \cos\theta_{ij}.$$
% Centering is what makes this identity work: the standard correlation formula
% subtracts the column means before taking inner products, but if the columns
% are {\em already} centered, that subtraction is a no-op, and what remains is
% the cosine-similarity formula on $\mathbb{R}^n$.
%
% Two structural points worth holding alongside the identity. (i) The columns
% live in $\mathbb{R}^n$ --- one entry per observation --- so the geometry of
% correlation is the geometry of the {\em observation space}, not the feature
% space. (ii) Correlation is normalized; covariance is not. The numerator
% $\langle \mathbf{x}_i, \mathbf{x}_j \rangle$ gives covariance (up to a factor
% of $n$ or $n-1$); dividing by the norms gives correlation.
%
% [Conceptual through-line: this is the bridge between Wk 5 (cosine similarity
% as the canonical geometric measure of alignment between vectors) and Wk 11
% (correlation as the canonical statistical measure of linear association).
% They are the same object. Whenever a student computes a correlation, they
% are computing a cosine; whenever they reason about correlated features,
% they are reasoning about angles in $\mathbb{R}^n$.]
%
% \textbf{Trick answer:} E (D5 naive-prior: drops the normalization entirely,
%   conflating correlation with covariance. Covariance --- not correlation ---
%   is the inner product (up to a $1/n$ or $1/(n-1)$ scaling); correlation
%   additionally divides by the norms. Catches students who half-remember
%   ``correlation has something to do with inner products'' without tracking
%   what makes it dimensionless.)
% \textbf{Trick answer:} C (D6 shape-implies-property: rows of $X$ are
%   observations in $\mathbb{R}^d$, not features. The angle between rows is
%   a geometric quantity, but it measures similarity between {\em observations},
%   not the correlation between {\em features}. The classic feature/observation
%   axis confusion in PCA, parallel to the $X^T X$ vs $X X^T$ trap in P32.)
% \textbf{Trick answer:} B (D5 naive-prior: imports an asymmetric projection
%   picture. The projection length $\langle \mathbf{x}_i, \mathbf{x}_j \rangle /
%   \norm{\mathbf{x}_j}$ is asymmetric in $i$ and $j$ --- but $\rho_{ij}$ is
%   symmetric. Catches students with a partial Wk 6 picture of orthogonal
%   projection who haven't connected it to the symmetric cosine.)
% \textbf{Trick answer:} D (D7 convention collision: imports the standardization
%   step of correlation PCA into the definition of $\rho_{ij}$. Standardization
%   --- rescaling each column to unit norm or unit sample variance --- is a
%   {\em preprocessing} choice that changes which PCA you run; it is not a
%   prerequisite for $\rho_{ij}$ to equal a cosine. The cosine identity holds
%   for centered columns regardless of their norms.)
\divider

{\bf PROBLEM 43:}
% WEEK = 7
% TOPICS = Eigenvalue characterization, singularity of A - lambda I,
%          nontriviality of the kernel
% DIFFICULTY = Core
%
Let $A$ be an $n \times n$ matrix and $\lambda \in \mathbb{R}$ a scalar. Which of the following is equivalent to ``$\lambda$ is an eigenvalue of $A$''?

\begin{enumerate}[(A)]
\item $\det(A) = \lambda^n$
\item $A\mathbf{v} = \lambda \mathbf{v}$ for every $\mathbf{v} \in \mathbb{R}^n$
\item $A - \lambda I$ is singular
\item $(A - \lambda I)\mathbf{v} = \mathbf{0}$ has only the trivial solution $\mathbf{v} = \mathbf{0}$
\item None of the above
\end{enumerate}

% \textbf{Correct Answer:} C
%
% \textbf{Explanation:} The chain of equivalences that licenses eigenvalue
% computation is: $\lambda$ is an eigenvalue $\Leftrightarrow$ there exists
% a \emph{nonzero} $\mathbf{v}$ with $(A - \lambda I)\mathbf{v} = \mathbf{0}$
% $\Leftrightarrow$ $A - \lambda I$ has a nontrivial kernel
% $\Leftrightarrow$ $A - \lambda I$ is singular $\Leftrightarrow$
% $\det(A - \lambda I) = 0$. The bridge from the geometric definition
% (an invariant scaling direction) to the computational recipe (root
% of the characteristic polynomial) runs through the singularity of
% $A - \lambda I$. That is what choice (B) names directly.
%
% Without this bridge, the characteristic polynomial feels like an
% unmotivated trick. With it, eigenvalues are exactly the scalars that
% ``collapse'' $A - \lambda I$ --- every other property follows.
%
% \textbf{Trick answer:} D (the exact \emph{negation} of the eigenvalue
%   condition. ``$(A - \lambda I)\mathbf{v} = \mathbf{0}$ has only the
%   trivial solution'' says $\ker(A - \lambda I) = \{\mathbf{0}\}$,
%   i.e., $A - \lambda I$ is \emph{invertible} --- which is precisely
%   the condition under which $\lambda$ is \emph{not} an eigenvalue.
%   Catches students who recognize that eigenvalues involve the
%   equation $(A - \lambda I)\mathbf{v} = \mathbf{0}$ but lose track
%   of the nontriviality requirement on $\mathbf{v}$. The closest
%   miss possible: right equation, wrong quantifier.)
% \textbf{Trick answer:} B (replaces ``some nonzero $\mathbf{v}$'' with
%   ``every $\mathbf{v}$.'' Universal quantification would mean
%   $A = \lambda I$, a vastly stronger condition than having $\lambda$
%   as an eigenvalue. Catches students who don't read quantifiers
%   carefully --- twin pair to (A) at the quantifier level.)
% \textbf{Trick answer:} A (uses the false ``determinant is
%   multiplicative across $\lambda I$'' rule: $\det(A - \lambda I)
%   \neq \det(A) - \lambda^n$, and $\det(A) = \lambda^n$ has no
%   general meaning. A student applying a phantom algebraic
%   identity lands here.)
% \textbf{Trick answer:} E (rewards the second-guesser who reads (B)
%   correctly but talks themselves out of it; the standard ``None''
%   trap when (B) is the genuine answer.)
\divider



\end{document}